%==========================================================================
% 基于动态图注意力网络与多智能体强化学习的低空交叉路口飞行器协同管控
% 期刊论文版式（依据《论文评审意见稿》逐条修订，v16）
% 作者：彭义森  指导教师：韩云祥  四川大学
% 编译方式：xelatex（两遍）
%==========================================================================
\documentclass[12pt,a4paper,UTF8]{ctexart}

%--------------------------------------------------------------------------
% 宏包
%--------------------------------------------------------------------------
\usepackage{geometry}
\usepackage{amsmath,amssymb,bm}
\usepackage{graphicx}
\usepackage{booktabs}
\usepackage{hyperref}
\usepackage{float}
\usepackage{enumitem}
\usepackage{multirow,array}
\usepackage{cite}
\usepackage{caption}
\usepackage{subcaption}
\usepackage{threeparttable}
\usepackage{algorithm}
\usepackage{algorithmic}
\usepackage{tikz}

\geometry{left=3cm,right=2.5cm,top=2.5cm,bottom=2.5cm}
\hypersetup{colorlinks=true,linkcolor=blue,citecolor=blue,urlcolor=blue}
\pagestyle{plain}

\renewcommand{\baselinestretch}{1.45}   % 适度行距，兼顾可读性与排版密度

\newcommand{\dyg}{DyGAT}

\begin{document}

%--------------------------------------------------------------------------
% 论文标题与作者（沿用v15版式）
%--------------------------------------------------------------------------
\title{\textbf{基于动态图注意力网络与多智能体强化学习的\\低空交叉路口飞行器协同管控}}
\author{彭义森\quad 指导老师：韩云祥\\[4pt]\small 四川大学}
\date{2026年7月}

\maketitle

%--------------------------------------------------------------------------
% 中文摘要
%--------------------------------------------------------------------------
\vspace{6pt}
\noindent\textbf{摘\quad 要：}城市低空交叉路口中多架电动垂直起降飞行器（eVTOL）的高密度安全协同管控是先进空中交通（AAM）面临的关键技术挑战。针对现有方法在飞行器间时空耦合建模与稀疏冲突样本利用两方面的不足，本文提出一种融合动态图注意力网络（DyGAT）的耦合多智能体强化学习（MARL）协同管控模型：以飞行器对之间的距离、速度差与航向差构造乘法门控注意力掩码，对学习注意力进行软性调制，将有限表征能力集中于决策敏感的交互对，同时软性抑制信息冗余的必然冲突对；并设计速度自适应安全距离、分层奖励函数，以及课程学习与Episode级冲突损失加权训练机制。在21架eVTOL交叉路口仿真场景中，模型与9种基线方法系统对比，无冲突完成率（CFCR）由最优基线MAPPO的$57.8\%$提升至$68.4\%$（$p<0.001$，Cohen's $d=5.27$），最小间隔违反率（MSVR）由$3.0\%$降至$2.4\%$（$p=0.029$，$d=1.69$），均达到统计显著；总冲突次数（TC）较MAPPO降低$10.3\%$，效应量达$d=0.93$，但受$n=5$样本量限制未达统计显著（$p=0.180$），属于方向性证据；消融实验表明DyGAT模块整体为贡献最大的组件，移除后总冲突次数增加$173.6\%$。结果表明，所提方法能够显著提升高密度低空交叉路口的飞行安全水平，为低空交通的智能化协同管控提供了一种可行途径。

\vspace{6pt}
\noindent\textbf{关键词：}低空交通管控；城市空中交通；多智能体强化学习；动态图注意力网络；乘法门控注意力掩码；冲突避免

\vspace{8pt}
\hrule
\vspace{10pt}

%--------------------------------------------------------------------------
% 英文标题与摘要（中文摘要之后，供期刊投稿使用）
%--------------------------------------------------------------------------
\begin{center}
{\Large\bfseries Collaborative Control of Aircraft at Low-Altitude Intersections\\ Based on Dynamic Graph Attention Networks and Multi-Agent Reinforcement Learning}
\end{center}
\vspace{4pt}
\noindent\textbf{Abstract:} The high-density safe and cooperative control of multiple electric vertical take-off and landing (eVTOL) aircraft at urban low-altitude intersections is a key technical challenge for advanced air mobility (AAM). To address the shortcomings of existing methods in modeling spatio-temporal coupling between aircraft and exploiting sparse conflict samples, this paper proposes a coupled multi-agent reinforcement learning (MARL) model that incorporates a dynamic graph attention network (DyGAT). A multiplicative gating attention mask is constructed from the pairwise distance, speed difference and heading difference, which softly modulates the learned attention so that limited representational capacity is concentrated on decision-critical interaction pairs while redundant inevitable-conflict pairs are softly suppressed. A speed-adaptive safety distance, a hierarchical reward function, and a training mechanism combining curriculum learning and episode-level conflict-loss weighting are further designed. In a simulated intersection scenario with 21 eVTOL aircraft, the proposed model is systematically compared against 9 baselines. The conflict-free completion rate (CFCR) is improved from $57.8\%$ (best baseline MAPPO) to $68.4\%$ ($p<0.001$, Cohen's $d=5.27$), and the minimum separation violation rate (MSVR) is reduced from $3.0\%$ to $2.4\%$ ($p=0.029$, $d=1.69$), both statistically significant; the total conflicts (TC) decrease by $10.3\%$ with a large effect size ($d=0.93$) but fail to reach significance at $n=5$ ($p=0.180$), which is treated as directional evidence. Ablation shows that the DyGAT module as a whole is the largest contributing component, whose removal increases the total number of conflicts by $173.6\%$. These results indicate that the proposed method can significantly improve flight safety in high-density low-altitude intersections and offers a feasible path toward intelligent cooperative control of low-altitude traffic.
\vspace{4pt}
\noindent\textbf{Keywords:} low-altitude traffic management; urban air mobility; multi-agent reinforcement learning; dynamic graph attention network; multiplicative gating attention mask; conflict avoidance
\vspace{8pt}
% \hrule
% \vspace{10pt}

%==========================================================================
% 1 引言
%==========================================================================
\section{引言}

\subsection{研究背景与意义}

近几年，低空经济被列入国家战略性新兴产业，电动垂直起降飞行器（Electric Vertical Take-Off and Landing，简称eVTOL）也随之从实验室原型快步走向商业示范。eVTOL是一种能像直升机一样垂直起降、又像固定翼飞机一样高效巡航的电动飞行器，不需要跑道，能在楼顶、停车场这类城市空间起降，因此被视为城市空中交通（Urban Air Mobility，UAM）\cite{ref45}的核心载运工具。亿航EH216-S已经在国内拿到了全球首张无人驾驶eVTOL型号合格证，Joby、Volocopter、Archer等国外企业也陆续进入适航审定阶段。美国联邦航空管理局（FAA）在《城市空中交通运行概念2.0》中给出了UAM从试点到高密度城市运营的演进路径，并提出协同运行区（Cooperative Areas）与可扩展交通管理（xTM）等运行框架\cite{ref1}；各类行业预测普遍认为，随着机队规模扩大，低空空域的运行架次将达到百万量级。

机队规模一旦达到这种量级，一个现实问题就会立刻浮出水面：如此多的飞行器挤在真高数百米以下的低空空域里，靠什么来保证它们互不碰撞、有序通行？传统民航空管靠管制员语音指挥，间隔以公里计，难以适应低空单位空域内数十架飞行器、运行状态在秒级尺度快速变化的密度水平。低空交通管理系统（UTM）因此被寄予厚望——它要求用自动化系统实时感知空域态势、自动做冲突探测与解脱。而在所有低空场景中，交叉路口又是冲突最集中、最棘手的地方：多股交通流在路口中心交汇，任何两架相邻方向飞来的飞行器，航向线都会在中心附近交叉，冲突几乎不可避免。如何让这些飞行器在保证安全间隔的前提下尽量不打乱节奏地通过路口，是先进空中交通（Advanced Air Mobility，AAM）落地前必须解决的关键问题。

从问题的本质看，多架飞行器在共享空域中的协同序贯决策存在强博弈耦合：每一架飞行器的最佳动作，都取决于其他飞行器接下来会怎么飞。它既要保证安全，又要兼顾效率，还要在决策周期只有几百毫秒的限制下实时求解。图神经网络适合建模飞行器间的结构化空间关系，注意力机制适合筛选关键交互信息，多智能体强化学习（MARL）适合序贯决策，如何把这三种能力有机融合，是智能交通与人工智能交叉领域的前沿问题。从工程角度看，本文的研究直接服务于UTM的核心功能模块：当机队规模达到数十架甚至上百架时，人工管制在信息处理带宽上已接近极限，基于规则与数学优化的传统方法又难以在计算复杂度与求解最优性之间取得平衡。一个能够实时、自主、协同地生成无冲突航迹的智能管控模型，是低空物流网络、城市空中交通示范运营乃至应急救援体系大规模落地的前提条件。

\subsection{本文主要贡献}

针对低空交叉路口的飞行器协同管控问题，本文的主要贡献概括如下：

（1）构建了面向低空交叉路口冲突规避的仿真环境（AirTrafficEnv）。该环境引入与速度相关的动态安全距离模型，使安全间隔随飞行器平均速度自适应调整；纳入最大加速度、最大减速度与最大转弯角等eVTOL物理运动约束，避免动作脱离真实飞行能力；并叠加高斯传感器噪声与随机风扰，使仿真条件尽量贴近真实低空运行。

（2）提出动态图注意力网络（DyGAT）。它在经典图注意力网络（GAT）基础上引入基于物理量的乘法门控注意力掩码：利用三个手工设计的指数核，把飞行器对之间的距离、速度差和航向差编码为一个$(0,1]$区间的软性掩码，以乘法方式作用于学习注意力，将有限表征能力集中于决策敏感的交互对。该机制不引入任何可学习参数，消融实验表明，去掉整个DyGAT模块后总冲突次数增加$173.6\%$。其中，将物理先验以零参数乘法掩码的形式注入注意力并非在一切领域都属首次（见第2.5节），本文将其应用于低空交叉路口eVTOL协同管控场景，并给出物理语义明确的三因子构造。

（3）设计课程学习与Episode级冲突损失加权相结合的训练机制。课程学习使模型从$N=5$架起步、逐步扩展至$N=21$架，降低初始学习难度；Episode级冲突加权（$\lambda=2.0$）对含冲突的Episode施加更大损失权重，在不引入off-policy数据、不破坏PPO on-policy框架的前提下，缓解稀疏冲突样本带来的训练困难。

（4）建立双层评价指标体系。在传统表层指标（累计奖励、总冲突次数、平均延误）之外引入深层安全指标（无冲突完成率、最小间隔违反率、平均冲突幅度），并据此揭示PCR（$99.1\%$）与CFCR（$68.4\%$）之间约30个百分点的系统性差距，对CTDE-MARL安全评估方法论提出了警示。

（5）与9种基线方法进行系统对比，并给出包含Welch $t$检验、Cohen's $d$效应量与统计功效分析的统计证据链，同时如实报告TC与MCAE两项指标在$n=5$下未达显著的结果，确保性能结论既统计可信、又如实反映小样本下的证据强度。

本文其余内容组织如下：第2节梳理相关工作并指出已有研究的不足；第3节将问题形式化为分布式马尔可夫决策过程并介绍仿真环境与奖励函数设计；第4节详细阐述DyGAT-MARL模型的结构与训练机制；第5节给出实验设置与结果分析；第6节进行讨论并总结全文。

%==========================================================================
% 2 相关工作
%==========================================================================
\section{相关工作}

围绕低空/空中交通的协同管控，学术界已经形成了“传统优化—工程安全—数据驱动”三条比较清晰的路线。本节先交代低空交通运行与安全间隔的基本概念，再沿着上述脉络梳理已有研究，并把本文用到的基础概念解释清楚，为后文铺垫。

\subsection{低空交通运行与安全间隔的基本概念}

在讨论具体方法之前，有必要先厘清几个贯穿全文的基础概念。

\textbf{低空空域与运行高度。}我国2023年发布的《国家空域基础分类方法》将空域划分为A、B、C、D、E、G、W七类，其中G类为非管制空域（B、C类空域以外真高300 m以下，W类除外），W类为G类空域内真高120 m以下、供微型/轻型/小型无人驾驶航空器运行的空域\cite{ref46}。行业通常将真高1000 m（含）以下、可延伸至3000 m的空域统称为低空空域\cite{ref46,ref48}，其中120--300 m被视为物流配送与城市空中交通的主要高度带，120 m以下主要为微型与轻型无人机的运行高度（依据《民用无人驾驶航空器运行安全管理规则》（CCAR-92）对微型无人机飞行高度的限制）\cite{ref47}。与民航高空航路（飞行高度通常在6000 m以上、航路之间间隔以公里计）相比，低空空域在垂直方向上的可用余量极其有限，多数eVTOL在巡航阶段采用固定的高度层，实际可用的垂直间隔资源很少——这就使得水平方向的协同避让成为低空冲突解脱的主要手段，也是本文将问题简化为二维平面进行建模的现实依据。

\textbf{安全间隔与冲突。}安全间隔（separation）是空管领域最基础的概念，指两架飞行器之间必须保持的最小水平或垂直距离。在传统民航中，雷达管制下的最小水平间隔通常为5海里（约9.3 km），该取值适用于区域雷达管制环境并视雷达覆盖与运行程序而定（ICAO Doc 4444）\cite{ref49}。对低空eVTOL而言，由于其速度低、机动性好、且机身体量远小于民航客机，安全间隔可以大幅缩小，但具体数值仍须综合考虑机体尺寸、旋翼下洗气流、传感器精度等因素。本文引入的“速度自适应安全距离”概念即是对这一思路的建模：间隔不再是固定常数，而是随飞行速度动态调整（详见第3.3.3节）。所谓“冲突”（conflict），是指两架飞行器之间的距离小于安全间隔这一状态；它与“碰撞”（collision）有本质区别——碰撞意味着已经发生物理接触，而冲突是可预测、可避免的潜在危险状态。空域管理的基本目标，就是通过及时有效的调控，把“冲突”消灭在发生之前，避免其升级为“碰撞”。

\textbf{低空交通管理系统（UTM）。}UTM是由美国联邦航空管理局（FAA）推动的面向无人机系统的交通管理框架，其核心思想是把传统依赖人工管制的空管体系，转变为数字化、自动化的服务体系：系统实时收集各飞行器的位置、航向、意图等信息，提供飞行计划申报、态势感知、冲突探测与解脱等自动化服务。当机队规模达到数十甚至上百架时，人工指挥在时间尺度和信息带宽上均无法胜任，这就决定了低空交通管控必须走向算法决策。本文研究的问题——交叉路口多机协同管控——即UTM中“冲突探测与解脱”这一核心服务在高冲突场景下的具体化。

\textbf{时间尺度。}低空飞行器的一个突出特点是决策周期极短。以本文场景为例，控制周期为0.5 s，一架以100 km/h飞行的eVTOL在0.5 s内即可移动约14 m——这意味着管控系统必须在亚秒级时间尺度内完成“感知—决策—执行”的闭环。这一严苛的实时性约束，贯穿了本文环境建模（第3章）与模型设计（第4章）的始终。

\subsection{传统空中交通流量管理方法}

空中交通流量管理（Air Traffic Flow Management，ATFM）是空管领域最早成体系的优化框架，核心思路是在已知空域容量的前提下，用整数规划、线性规划等数学手段，对航班的地面等待、改航和调速做集中式决策。Menon等\cite{ref2}把航空交通流抽象成连续介质模型，用偏微分方程描述交通流密度随时间的演化，再求解最优流量控制策略。这套方法在高空航路网络里很有效，但前提是“空域状态全局可知、决策由单一中心节点完成”。一旦飞行器数量上升到几十上百架，状态空间和约束规模就会指数膨胀，实时求解基本不可能。这种“集中式优化难以扩展”的固有短板，也是后来数据驱动方法切入的地方。

在地面交通信号控制领域，自适应信号控制的研究为空中交通管控提供了重要的方法论借鉴。Aslani等\cite{ref12}把actor-critic强化学习算法用到真实交通网络的信号控制上，验证了它在多种交通扰动下的鲁棒性；国内学者也较早探索了基于Q-learning的交叉口信号控制\cite{ref16,ref19,ref20}，并逐步过渡到深度强化学习\cite{ref15,ref17,ref18,ref22}。这些工作表明，强化学习在处理多路口、多时段、强耦合的交通控制问题方面具有优势，这也为本文将其思想推广到空中交通场景提供了依据。

\subsection{基于工程约束的冲突解脱方法}

与宏观流量管理相对，工程安全方法着眼于微观层面的冲突解脱。它们的共同点是拿明确的物理或数学约束做安全保证，但各有各的适用边界。

\textbf{模型预测控制（MPC）。}Richards与How\cite{ref3}把飞行器航迹规划与冲突避免建模成混合整数线性规划，在滚动时域里反复求解最优控制序列，每步只执行当前时域的第一个控制量、随后滚动重算。它的优点是约束表达清晰、最优性可度量，缺点是计算量随预测时域和飞行器数量快速增长，在高密度低空场景下实时性难以保证。

\textbf{速度障碍法（VO）与最优互惠避碰（ORCA）。}速度障碍法的思想很直观：如果两架飞行器保持当前相对速度不变就会相撞，那么所有“会导致相撞的相对速度”构成一个障碍集合，避碰就是选择一个不落在集合里的速度。具体而言，设飞行器$A$、$B$的位置与速度分别为$p_A,v_A$与$p_B,v_B$，将$B$视为以$A$为参考的一个圆盘，若$A$的相对速度$v_A-v_B$指向该圆盘，则$A$在保持当前相对速度的情况下必然碰撞——所有这样的相对速度构成的集合，即为$A$对$B$的速度障碍$VO^B_A$。van den Berg等\cite{ref4}在此基础上提出ORCA，引入“互惠”假设——双方各承担一半避让责任，于是每架飞行器只需在其可行速度空间中选择与期望速度偏差最小、且位于避碰约束半平面内的速度，通过线性规划即可高效求解，实现了分布式的平滑避碰。ORCA已广泛用于机器人与无人机导航，但它依赖“所有智能体对称地分担避让责任”这一假设，一旦智能体行为异构或目标相互冲突，避让质量就会明显下降。

\textbf{控制屏障函数（CBF）。}Ames等\cite{ref5}系统总结了CBF理论。它的核心思想是设计一个满足前向不变性条件的“安全屏障”函数，把安全约束编码成在线可解的二次规划约束，在尽量不干扰名义控制器的前提下保证系统始终待在安全集合内。具体而言，设安全集合为$C=\{x:h(x)\geq0\}$，通过构造满足前向不变性条件的屏障函数$h$，将安全约束$L_f h(x)+L_g h(x)u+\alpha(h(x))\geq0$作为约束加入名义控制器，即可保证系统状态始终停留在安全集合内。CBF的安全保证是严格的，但对系统模型精度要求很高，而且多架飞行器共享空域时，各个安全集合的交集会迅速收缩，可行解很容易退化。

总的来说，工程安全方法在形式化安全保证方面具有不可替代的优势，但面对高密度、强动态的交互场景都不同程度地存在局限，这为学习方法留出了空间。

\subsection{强化学习与多智能体强化学习在交通管控中的应用}

强化学习（Reinforcement Learning，RL）让智能体通过与环境的试错交互来学习策略，不需要精确建模系统动力学。它背后的数学模型是马尔可夫决策过程（Markov Decision Process，MDP）：用状态、动作、转移概率、奖励和折扣因子描述“智能体在什么状态、做什么动作、获得多少回报、环境怎么变”这一整套循环。MDP满足马尔可夫性，即下一状态仅依赖于当前状态与动作、与更早的历史无关，这使得智能体无需记忆完整历史即可做出最优决策。智能体的目标，是学到一个能最大化长期累计回报（即各步奖励按折扣因子加权求和）的策略。

在MDP框架中，衡量状态与动作好坏的核心工具是价值函数。状态价值函数$V^\pi(s)$表示从状态$s$出发、按照策略$\pi$继续执行所能获得的期望累计回报；动作价值函数（Q函数）$Q^\pi(s,a)$则表示从状态$s$执行动作$a$后再按策略$\pi$执行的期望回报。二者满足著名的贝尔曼方程（Bellman equation），它揭示了相邻时刻价值之间的递推关系。这一方程几乎是所有强化学习算法的理论基石——无论是基于价值的方法（如DQN）通过估计Q函数间接得到策略，还是基于策略的方法（如策略梯度、PPO）直接对策略参数求梯度更新，其核心更新规则都可追溯到贝尔曼方程。当动作空间为连续值（如本文的加速度指令）时，对Q函数求最大值本身就是一个困难的优化问题，这促使我们选择直接优化策略的策略梯度类方法，这也是后文采用PPO的深层原因。

在单智能体层面，深度Q网络（DQN）\cite{ref40}及其变体被广泛用于冲突规避和信号控制\cite{ref12,ref15,ref16,ref17,ref18,ref19,ref20,ref21,ref22}。江波等\cite{ref21}将深度Q学习应用于机场道口的协同调速，验证了RL在航空地面场景的可行性。但单智能体RL存在一个固有局限：它把其他飞行器当作环境的一部分，无法刻画它们之间的博弈耦合关系。当每架飞行器都只优化自身收益时，联合行为往往偏离系统最优，甚至引发冲突的连锁反应。

多智能体强化学习（Multi-Agent Reinforcement Learning，MARL）正是为解决这个问题而生的。多智能体场景的核心困难在于非平稳性：其他智能体的策略也在不断更新，使得每个智能体面对的环境始终处于动态变化之中，破坏了单智能体强化学习所依赖的马尔可夫性。另一大困难是信用分配，即如何把团队共同获得的收益，合理归因到每个智能体的具体行为上。为应对这些难题，学术界发展出集中训练分布式执行（Centralized Training with Decentralized Execution，CTDE）范式：训练阶段共享全局信息，把智能体间的耦合关系显式建模，使训练信号更稳定；执行阶段每个智能体仅依据自身局部观测独立决策，以满足实时性与通信约束。围绕CTDE，Tumer与Agogino\cite{ref6,ref7,ref14,ref24}较早验证了分布式多智能体学习在ATM中的有效性；Spatharis等\cite{ref11}、Ghavamzadeh等\cite{ref23}设计了分层MARL方案；Brittain与Wei\cite{ref8}在高密度航路扇区验证了深度MARL的自主间隔保障；Ghosh等\cite{ref9}用深度集成MARL提升了管制决策的稳健性，Xie等\cite{ref29}则将强化学习用于UAM与密集低空交通的流量管理。算法层面，MADDPG（Multi-Agent Deep Deterministic Policy Gradient）\cite{ref34}采用集中式Critic与分布式Actor的结构，训练时Critic借助全局信息评估联合动作的价值，执行时每个Actor仅依据自身观测输出动作；COMA（Counterfactual Multi-Agent Policy Gradients）\cite{ref35}通过反事实基线——比较“智能体实际动作带来的回报”与“若它换作其他动作的期望回报”——衡量每个智能体对团队收益的真实边际贡献，从而解决信用分配问题；QMIX\cite{ref33}通过单调混合网络将全局Q值分解为个体Q值的单调组合，要求全局Q值随任一个体Q值单调不减，以保证分解与全局最优一致；MAPPO（Multi-Agent PPO）\cite{ref36}则把单智能体PPO的稳定训练特性推广到多智能体场景，是目前合作博弈任务中的强基线。国内方面，翟子洋等\cite{ref13}系统梳理了大规模智慧交通信号控制中强化学习的技术演进。

\subsection{图神经网络、注意力机制与先验注入}

飞行器之间的关系天然适合用“图”来表达：每架飞行器是节点，飞行器对之间的交互是边。形式化地说，图记作$\mathcal{G}=(\mathcal{V},\mathcal{E})$，其中$\mathcal{V}$为节点集合、$\mathcal{E}$为边集合，邻接矩阵$\bm{A}\in\mathbb{R}^{N\times N}$刻画节点间的连接关系（$A_{ij}=1$表示节点$i$与$j$相连，$A_{ij}=0$表示不相连）；当边带有权重时，$A_{ij}$可取连续值——这一“连续边权”即后文DyGAT的关键所在。图神经网络（GNN）就是专门在这种结构上做特征学习的模型。Kipf与Welling\cite{ref32}提出的图卷积网络（GCN）通过邻居聚合更新节点特征，但聚合权重只由图的拓扑结构（即节点度数，也就是邻居数量）决定，不随节点特征变化，因此难以区分重要邻居与次要邻居。Veli\v{c}kovi\'c等\cite{ref31}提出的图注意力网络（GAT）解决了这个问题——它让网络自己学习“哪些邻居更值得关注”，通过注意力权重自适应地聚合信息。

注意力机制背后还有一个与本文密切相关的重要思想：门控。长短期记忆网络（LSTM）\cite{ref37}之所以能记住长序列信息，靠的就是遗忘门、输入门、输出门这一套门控结构，通过控制“历史信息保留多少、新信息写入多少”来调节信息流。这类门控信号都是通过训练学出来的。本文尝试一个相反的思路：不用学习，直接让物理量充当门控信号，把飞行中的基本安全常识注入注意力计算，这即后文DyGAT的核心思想。

将先验信息注入注意力或图结构并非本文在任何意义上的首创：机器人避障与自动驾驶领域已有面向危险度的注意力（danger-aware attention）、物理信息图神经网络（physics-informed GNN）等相近思想\cite{ref55}，多智能体领域也有基于智能体关系的注意力掩码。本文的差异化体现在三个方面：其一，先验取自飞行安全常识（距离、速度差、航向差），物理语义明确，且与低空eVTOL的具体运行约束（局部感知、亚秒级决策）直接对应；其二，门控完全零参数，由状态量直接计算，不需要额外训练；其三，门控与课程学习、冲突损失加权构成完整的训练机制。因此，本文对“首次”的表述严格限定在“低空交叉路口eVTOL协同管控场景下的乘法门控注意力机制”这一具体语境，不对一般意义上的物理先验注入主张首创。

在低空无人机交通领域，Li等\cite{ref10}提出的MS-GRL方法用多尺度GAT增强强化学习，解决了100架UAV的密集冲突消解问题，是图增强MARL的代表性工作。但它和大多数GNN方法一样，图结构基于固定距离阈值构建——只要两机距离在阈值以内就一律视为邻居。其局限在于：它区分不了“距离相同但风险不同”的飞行器对。两对飞行器当前欧氏距离完全一样，一对相向高速靠近、另一对同向等速，冲突风险能差出几十倍，固定阈值图却把它们一视同仁。这是本文DyGAT要解决的问题。

\subsection{与代表性相关工作的系统性对比}

为更直观地展示本文方法与代表性工作的差异，表\ref{tab:related_work}从五个维度对方法进行了横向对比：时空耦合建模能力、图结构的动态性、稀疏奖励的处理方式、物理约束的完整程度以及扰动建模的丰富程度。其中，“完整”的物理约束与“风扰+噪声”扰动建模仍是对真实低空环境的简化——Dryden标准风谱、城市峡谷效应等更复杂的因素尚未纳入，这一简化在第3.3.3节有专门讨论，其影响留待后续以标准风谱模型做对照实验检验。

\begin{table}[H]
    \centering
    \caption{本文方法与代表性相关工作的系统性对比}
    \label{tab:related_work}
    \resizebox{0.95\textwidth}{!}{%
    \begin{tabular}{|c|c|c|c|c|c|}
        \hline
        \textbf{方法} & \textbf{时空耦合建模} & \textbf{图结构类型} & \textbf{稀疏奖励处理} & \textbf{物理约束} & \textbf{扰动建模} \\
        \hline
        Brittain \& Wei \cite{ref8} & LSTM+FC & 静态全连接 & $\times$ & 部分 & $\times$ \\
        \hline
        Ghosh et al.\ \cite{ref9} & Geo-hash密度图 & 网格化 & $\times$ & 部分 & $\times$ \\
        \hline
        Spatharis et al.\ \cite{ref11} & 预定义关系矩阵 & 层级分组 & $\times$ & $\times$ & $\times$ \\
        \hline
        MS-GRL \cite{ref10} & 多尺度GAT & 固定阈值图 & $\times$ & 部分 & $\times$ \\
        \hline
        QMIX/MAPPO基线 & FC/RNN & 全连接 & $\times$ & 部分 & $\times$ \\
        \hline
        \textbf{本文方法} & \textbf{DyGAT（空间+时序+动态边权重）} & \textbf{动态图} & \textbf{课程学习+冲突轨迹加权} & \textbf{完整} & \textbf{风扰+噪声} \\
        \hline
    \end{tabular}%
    }
\end{table}

从表\ref{tab:related_work}可以看出：已有方法或使用全连接/静态图建模智能体关系，或使用固定阈值构建图结构，在处理稀疏奖励与物理约束方面普遍留有空白；本文方法则在时空耦合建模、动态图结构、稀疏样本利用与物理约束注入四个维度上同时做出了针对性改进，形成与已有工作相区别的定位。

\subsection{现有研究的不足与本文切入点}

综合上面几节，已有研究在三个维度上还存在明显缺口：

（1）\textbf{时空耦合建模不足。}多数方法用全连接网络或静态图结构\cite{ref8,ref10,ref11}，区分不了“距离相同但相对运动状态不同”的飞行器对，在高密度动态场景下决策质量受限。

（2）\textbf{稀疏冲突样本利用不足。}冲突是低概率事件，在训练数据里占比通常不到$0.1\%$\cite{ref12,ref13}，随机采样下这些宝贵的危险样本很容易被大量正常样本淹没，难以提升模型对罕见危险的敏感性。

（3）\textbf{环境不确定性建模不足。}多数研究忽略物理运动约束和动态安全距离\cite{ref11}，风扰和噪声处理得也过于简单。安全间隔如果是个固定常数，就反映不出“飞得越快、越要多留余地”这个基本规律。

这三个缺口共同构成本文的切入点：时空耦合建模不足对应DyGAT的物理门控机制，稀疏样本利用不足对应课程学习与冲突损失加权，环境不确定性建模不足对应速度自适应安全距离与扰动建模。

%==========================================================================
% 3 问题建模与管控模型设计
%==========================================================================
\section{问题建模与管控模型设计}

\subsection{问题描述与场景设定}

本文聚焦城市低空十字交叉路口这一典型高冲突场景。如图\ref{fig:scenario}所示，$N$架eVTOL均匀分布在以交叉中心为圆心、半径$r_{\text{init}}$的圆周上，各自从不同方位径直飞向中心。飞行器的航向线天然在中心区域附近相互交叉——任何两架相邻或相对方位出发的飞行器都会在路口附近相遇，冲突威胁持续存在。飞行器需要在保持安全间隔的前提下，以尽可能小的延迟通过交叉区域。

\begin{figure}[H]
    \centering
    \begin{minipage}[c]{0.7\textwidth}
    \centering
    \begin{tikzpicture}[scale=0.95]
        \draw[->,thick] (-5.2,0) -- (5.2,0);
        \draw[->,thick] (0,-5.2) -- (0,5.2);
        \draw[dashed] circle (3.2);
        \foreach \ang in {0,60,120,180,240,300} {
            \fill[blue!70] ({3.4*cos(\ang)},{3.4*sin(\ang)}) circle (0.12);
        }
        \foreach \ang in {0,60,120,180,240,300} {
            \draw[->,blue!70,thick] ({3.2*cos(\ang)},{3.2*sin(\ang)}) -- ({1.1*cos(\ang)},{1.1*sin(\ang)});
        }
        \draw[red,very thick] (1.0,0.35) circle (0.6);
        \node[below right] at (1.7,1.05) {冲突区};
        \node[above] at (0,5.3) {交叉路口中心};
        \node[below] at (0,-5.7) {初始位置圆周};
    \end{tikzpicture}
    \end{minipage}
    \caption{低空十字交叉路口场景示意图。蓝色圆点为均匀分布于圆周上的eVTOL，箭头指向交叉中心，红色圆环为高冲突区域。}
    \label{fig:scenario}
\end{figure}

这个场景虽然抽象，但抓住了低空交通冲突的三个本质特征。第一是\textbf{对称冲突}：所有飞行器以相同优先级争夺中心区域，不存在天然的“让行次序”。第二是\textbf{密集交互}：$N$架飞行器最多构成$N(N-1)/2$个交互对，$N=21$时就是210对，远超人工调度的能力。第三是\textbf{时空耦合}：冲突风险不仅看当前距离，还看相对速度和相对航向——同样是相距100 m，相向而行和同向而行的风险天差地别。

为保证问题可解，本文作如下合理假设：所有飞行器在同一高度层运行（二维平面问题）；各eVTOL具有相同的动力学参数（同构智能体）；每架飞行器只能感知500 m半径内的邻居（局部可观测性），全局状态仅在训练阶段可用；训练阶段通信无延迟、无丢包。这些假设与真实eVTOL的机载感知能力基本一致，但同时也界定了方法的适用范围。这里需要明确两点简化边界。其一，二维平面假设忽略了三维高度层在低空交通中的潜在利用价值——真实低空中飞行器完全可以通过改变高度层来错开冲突，这是低空交通一种重要的安全冗余手段，也是本文暂未建模之处；其二，理想通信假设暂未考虑真实UTM部署中50--200 ms的通信延迟与可能的丢包，其对集中训练效果的影响有待后续检验。

\subsection{分布式马尔可夫决策过程形式化}

把上述问题写成分布式马尔可夫决策过程（Dec-MDP）\cite{ref41}，即六元组$\langle \mathcal{N}, \allowbreak \mathcal{S}, \allowbreak \mathcal{A}, \allowbreak \mathcal{P}, \allowbreak \mathcal{R}, \allowbreak \gamma \rangle$：

\begin{itemize}
    \item \textbf{智能体集合} $\mathcal{N} = \{1, 2, \ldots, N\}$。训练阶段$N$由课程学习从5逐步增至21，评估阶段固定为21。
    \item \textbf{全局状态} $\mathcal{S}$：$S = (s_1, \ldots, s_N)$，其中$s_i = [x_i, y_i, v_i, v_i^{\text{target}}, d_i^{\text{center}}]$为飞行器$i$的5维状态向量，含位置、当前速度、目标速度与到中心的距离。这个五维向量是形式化层面的状态定义；策略网络实际输入的是扩展后的局部观测$o_i$（第3.3.1节），全局状态$S$在CTDE训练阶段供GAT-Mixer计算全局价值函数使用（GAT-Mixer使用其位置与速度分量，见第4.4节）。
    \item \textbf{动作} $\mathcal{A}$：每架飞行器的动作为连续加速度$a_i$（$\mathrm{m/s^2}$），受运动学约束、风扰与噪声共同影响。
    \item \textbf{转移概率} $\mathcal{P}$：由第3.3节的运动学方程与不确定性模型隐式确定，不显式建模。
    \item \textbf{奖励函数} $\mathcal{R}$：团队奖励结构$R = \frac{1}{N}\sum_i r_i$，另有完成奖励（第3.4节）。
    \item \textbf{折扣因子} $\gamma = 0.99$，用于权衡近期奖励与远期奖励的相对重要程度，$\gamma$越接近1，模型越重视长期收益。
\end{itemize}

在进入具体设计之前，先说明三点。其一，本文所称“耦合”，指各飞行器的决策通过共享空域与团队奖励相互耦合：任何一架飞行器的动作都会通过改变全局态势而影响其他飞行器的决策，这种“你中有我”的相互依赖即多智能体协同区别于多架独立智能体并行优化的核心。其二，Dec-MDP在一般意义上是难以精确求解的。Bernstein等\cite{ref41}已经证明，在无通信约束下Dec-MDP的计算复杂度为NEXP完全（NEXP-complete）。直观地理解，每个智能体只能基于自己的局部观测决策，却要共同驱动一个全局状态的演化，这种“局部决策、全局耦合”的结构使得暴力搜索在规模上完全不可行。这也解释了为什么学术界普遍采用“集中训练、分布式执行”（CTDE）这一折中范式：训练阶段让智能体共享全局信息以简化学习，执行阶段再退化为纯局部决策以匹配现实约束。本文模型同样遵循这一范式，其具体体现是：训练时由GAT-Mixer聚合全局状态计算$Q_{\text{tot}}$，执行时各智能体仅依据自身观测输出动作。其三，关于团队平均奖励：如果每架飞行器只优化自己的安全，很容易出现“把风险甩给邻居”的机会主义——我减速让你撞上，只要我不撞就行。把个体奖励平均成团队奖励后，每一架的行为都会影响全局收益，客观上激励飞行器采取有利于整体交通流的策略。

\subsection{仿真环境设计}

仿真环境（AirTrafficEnv）以Python面向对象方式实现，是全部实验的基础平台。下面按观测空间、动作空间、不确定性建模三部分展开。

\subsubsection{观测空间设计}

观测空间分三个层次，层层递进。第一层是形式化状态$s_i\in\mathbb{R}^5$，即上面定义的五维向量。第二层是局部观测$o_i$。$o_i$由两部分构成：一是自身状态$s_i$；二是邻域信息——以$i$为圆心、半径500 m内的所有邻居$j$（不含自身，邻居集合记为$\mathcal{N}_i$）在最近$T=10$步（对应5 s）内的历史特征序列。对每个邻居$j$，其单步特征取相对位置$(\Delta x_{ij}, \Delta y_{ij})$、邻居速度$v_j$与邻居航向角$\theta_j$（共4维），其中航向角定义为相对交叉中心的方向角，与门控公式中的航向差语义一致。由于$\mathcal{N}_i$的大小随场景动态变化（$N=21$时平均约6），模型以“自身节点+各邻居节点”构造局部子图，节点特征为各自的LSTM时序嵌入（64维），节点数随邻域大小自适应变化，无需定长填充。第三层是编码特征，局部观测先经64维LSTM编码轨迹，再经256维DyGAT聚合空间信息，得到蕴含时空耦合信息的256维特征$h_i^{\text{coupled}}$：
\begin{equation}
    o_i \xrightarrow{\text{LSTM}_{64}} h_i^{\text{traj}} \xrightarrow{\text{DyGAT}_{256}} h_i^{\text{coupled}} \rightarrow \text{策略/Q值输出}.
    \label{eq:encoding_chain}
\end{equation}
其中$h_i^{\text{coupled}}$为DyGAT输出中对应自身节点的表示，与自身位置、速度、目标速度拼接后经决策头输出动作与Q值。观测假设：邻域内邻居的速度与航向经由V2V广播或地面UTM服务获取，这一假设与3.1节的理想通信假设一致；观测中不包含邻居加速度，因为加速度在现实场景中难以直接观测。

\subsubsection{动作空间与物理运动约束}

动作空间为连续加速度$a_i\in[-4.0, 3.0]$ $\mathrm{m/s^2}$，正值加速、负值减速。速度更新式为：
\begin{equation}
    v_i^{t+1} = \text{clip}\!\left(v_i^{t} + a_i \cdot \Delta t \cdot 3.6,\; 60,\; 140\right),
    \label{eq:velocity_update}
\end{equation}
其中$\Delta t=0.5$ s为控制周期，系数3.6用于把$\mathrm{m/s}$换算成$\mathrm{km/h}$，速度限制在60--140 $\mathrm{km/h}$。飞行器在二维平面内的运动由位置、速度与航向刻画。航向角$\theta_i$定义为速度方向与$x$轴正方向的夹角，速度分解为$v_i[\cos\theta_i,\ \sin\theta_i]^{\top}$。位置更新式为：
\begin{equation}
    x_i^{t+1} = x_i^{t} + v_i^{t+1}\,\frac{1000}{3600}\,\cos\theta_i^{t}\,\Delta t,\qquad
    y_i^{t+1} = y_i^{t} + v_i^{t+1}\,\frac{1000}{3600}\,\sin\theta_i^{t}\,\Delta t,
    \label{eq:position_update}
\end{equation}
其中$v_i$以$\mathrm{km/h}$计，$\times1000/3600$换算为$\mathrm{m/s}$。航向更新遵循“单步向交叉中心方向转向不超过$\pi/6$（30$^{\circ}$）”的限制：
\begin{equation}
    \theta_i^{t+1} = \theta_i^{t} + \text{clip}\!\left(\Delta\psi_i^{t},\; -\pi/6,\; \pi/6\right),
    \label{eq:heading_update}
\end{equation}
其中$\Delta\psi_i^{t}$为当前航向与指向交叉中心方向之间的带符号最小转向角。式(\ref{eq:heading_update})保证了表\ref{tab:env}中“最大转弯角$\pi/6$”这一约束，防止飞行器“瞬间掉头”这类不真实的运动。本文动作空间只包含纵向加速度$a_i$，航向由上述转向约束自动调节、始终朝交叉中心机动；因此观测与门控公式中使用的航向角$\theta_i$直接由运动状态导出，避免了把航向角也作为独立动作带来的控制维度膨胀。

速度范围参考了当前主流eVTOL的巡航速度；加速度取3.0、减速度取4.0 $\mathrm{m/s^2}$，减速度大于加速度是安全设计惯例。这里需要诚实说明：主流eVTOL厂商（Joby、Volocopter、Lilium、Archer等）都没有公开最大水平加速度的精确规格，学术界建模一般假设约2.0 $\mathrm{m/s^2}$\cite{ref25}。本文取值高于该文献假设，理由有三：(i)该文献针对早期eVTOL概念，近两年公布的新一代机型在动力冗余与爬升性能上普遍较早期概念有所提升；(ii)本文减速度大于加速度，使避让动作以减速为主，属于安全设计惯例；(iii)环境对航向变化施加$\pi/6$的独立约束，纵向加速度与转向能力解耦。仍需指出，在厂商公开精确规格之前，该取值属于工程估计；5.8节的参数敏感性分析未覆盖最大加速度，其影响留待后续实验检验。

\subsubsection{不确定性建模与动态安全距离}

真实低空环境存在两类不可回避的不确定性：传感器噪声与风扰。本文在观测通道叠加$\mathcal{N}(0, 0.5^2)$ m的高斯噪声，模拟GPS/视觉定位的典型误差；每架飞行器每步受到方向均匀随机、强度0.1 $\mathrm{m/s}$的风扰。对风扰的建模作如下简化：真实低空风场通常用Dryden（MIL-F-8785C）\cite{ref26}或Kaimal（IEC 61400-1）\cite{ref44}谱模型描述，具有空间相关性和时间连续性，本文采用的均匀随机风扰是简化版本，更精细的风场建模对方法性能的影响，留待后续以标准谱模型做对照实验检验。

安全间隔同样不该是个固定常数。直观上讲，飞得越快，单位时间内接近的距离越大，需要留的余量就该越多。据此定义速度自适应安全距离：
\begin{equation}
    d_{\text{safety}} = d_{\text{base}} + \bar{v} \times 0.5,
    \label{eq:safety_distance}
\end{equation}
其中$d_{\text{base}}=50$ m为基本安全间隔，$\bar{v}$为当前时刻场景内飞行器的平均速度，单位统一为$\mathrm{m/s}$，$\bar{v}\times0.5$可理解为“额外0.5 s前瞻时间内扫过的距离”。以$\bar{v}=70$ $\mathrm{km/h}$（约19.4 $\mathrm{m/s}$）为例，$d_{\text{safety}}\approx60$ m；加速到140 $\mathrm{km/h}$（约38.9 $\mathrm{m/s}$）时约69 m。低速不保守、高速有保护，这与空管实践中“速度越高、间隔越大”的原则完全一致。

为说明$d_{\text{base}}=50$ m取值的依据，可作如下对标。传统民航雷达管制下的最小水平间隔一般为5海里（约9.3 km，ICAO Doc 4444）\cite{ref49}，该间隔远大于eVTOL所需的机间距离；针对无人机的探测与避让（DAA）要求，RTCA DO-365给出了相应的系统性能标准\cite{ref51}，NASA提出的well clear包络约为水平数十英尺到数十米的量级；JARUS SORA则将安全评估与运行场景、地面人口密度等挂钩\cite{ref50}。综合机身尺寸（eVTOL翼展通常在10 m量级）、本环境的传感器定位误差（噪声标准差0.5 m）、决策反应时间（一个控制周期0.5 s）以及旋翼下洗等气动干扰，50 m的基本间隔在“远大于机身尺度、且能容纳传感器与控制延迟造成的偏差”的意义上是工程可接受的下限量级；$\bar{v}\times0.5$项进一步随速度线性增大间隔，与空管原则一致。

\subsection{分层奖励函数设计}

奖励函数是强化学习中连接行为与目标的桥梁，其设计质量直接影响学习效果。本文遵循三条原则：奖励信号要密集，以保证模型能频繁获得学习反馈；惩罚幅度要随危险程度增大，即越接近危险越重；各分量以自身表达式的尺度形成隐式权重，使安全、效率与协同能够权衡。据此构造四部分奖励：
\begin{equation}
    r_i = r_{\text{safe}} + r_{\text{speed}} + r_{\text{coop}} + r_{\text{eff}}.
    \label{eq:reward_total}
\end{equation}

式(\ref{eq:reward_total})的四个分量并未显式设置权重系数，其相对权重由各分量表达式自带的尺度决定：安全奖励的惩罚幅度远大于其他分量（危险区可达$-50e^{2}\approx-369$，而速度、协同、效率分量通常不超过$\pm20$），从而保证安全在综合奖励中占据主导地位；速度、协同与效率分量幅度接近，形成效率与协同的均衡。这一“以分量尺度为隐式权重”的设计使安全成为约束性目标、效率与协同成为优化性目标，符合“先安全后效率”的管控优先级。若需调整权衡，可显式引入权重系数，属于超参数调优的范畴。

\textbf{安全奖励}采用三段式分级结构：
\begin{equation}
    r_{\text{safe}} = \begin{cases}
        -50 \exp\!\left(2\frac{d_{\text{safety}}-d_{\min}}{d_{\text{safety}}}\right), & d_{\min}<d_{\text{safety}} \text{（危险区）}\\
        -10 \frac{1.2d_{\text{safety}}-d_{\min}}{0.2d_{\text{safety}}}, & d_{\text{safety}}\leq d_{\min}<1.2d_{\text{safety}} \text{（缓冲区）}\\
        +5, & d_{\min}\geq1.2d_{\text{safety}} \text{（安全区）}
    \end{cases}
    \label{eq:reward_safe}
\end{equation}
其中$d_{\min}$为飞行器$i$到所有邻居的最小距离。三个区域的物理含义很清晰：已经侵入安全间隔（危险区）时采用指数惩罚，距离越近惩罚越重，引导模型尽快拉开距离；处于警示缓冲带时采用线性惩罚，鼓励多留余量；保持充足间隔则给予正奖励。该分段函数在$d_{\min}=d_{\text{safety}}$与$d_{\min}=1.2d_{\text{safety}}$两处边界并不连续：危险区端值为$-50$，缓冲区起点为$-10$，存在40的阶跃；缓冲区末端为0，而安全区为$+5$，亦有跳跃。因此“避免了梯度不连续”这一表述并不严谨。实际上，RL对奖励值的不连续并不敏感——价值函数经网络平滑后仍可正常学习——该分段设计真正的作用是以清晰的分档信号把“接近危险”的程度传递给学习过程。此处修正表述：分段结构以指数+线性+正奖励的三档信号分级刻画危险程度。

\textbf{速度奖励}引导飞行器把速度维持在理想区间$[65,90]$ $\mathrm{km/h}$（与目标速度$70\pm10$ $\mathrm{km/h}$一致）：
\begin{equation}
    r_{\text{speed}} = \begin{cases}
        -50\frac{v_i-90}{50}, & v_i>90 \text{（超速）}\\
        -30\frac{65-v_i}{25}, & v_i<65 \text{（过慢）}\\
        +20, & 65\leq v_i\leq 90 \text{（理想）}
    \end{cases}
    \label{eq:reward_speed}
\end{equation}
超速会放大冲突风险，过慢会拖累交通流，只有落在理想区间才能兼顾安全与效率。

\textbf{协同奖励}刻画与邻居的速度一致性：
\begin{equation}
    r_{\text{coop}} = -0.2\left|v_i - \bar{v}_{\text{neighbor}}\right|,
    \label{eq:reward_coop}
\end{equation}
其中$\bar{v}_{\text{neighbor}}$为邻域内邻居的平均速度。这项奖励意在抑制“个别飞行器突进、整体节奏被打乱”的现象，在路口场景里有助于形成平滑的交通流。

\textbf{效率奖励}引导飞行器逼近各自的目标速度：
\begin{equation}
    r_{\text{eff}} = -0.5\left|v_i - v_i^{\text{target}}\right|,
    \label{eq:reward_eff}
\end{equation}
保证飞行器在规避冲突的同时仍以接近任务速度的方式飞行，避免“为了安全而过度保守”。

全局奖励为个体奖励的平均，$R=\frac{1}{N}\sum_i r_i$；当所有飞行器都通过交叉中心（距中心小于10 m）时，额外奖励+200。这个完成奖励属于稀疏奖励——只在Episode结束时出现一次——但幅度足够大，能让模型把“完成任务”这一长期目标纳入优化考量，与高密度的即时奖励形成互补。

\subsection{冲突检测与安全指标定义}

为对模型的安全性能进行细粒度评估，本节给出冲突检测的判定准则与严重程度分级方法。这些定义贯穿第5章的实验分析。

\textbf{冲突判定准则。}在仿真的每个控制周期，对任意飞行器对$(i,j)$计算欧氏距离$d_{ij}$，若$d_{ij}<d_{\text{safety}}$，则判定二者之间发生一次冲突。这里的安全距离采用第3.3.3节的速度自适应定义，因此“是否冲突”的判定标准会随飞行器的平均速度动态变化——这一设计使冲突检测与“飞得越快、间隔越该越大”的物理规律保持一致。此外，为便于分析模型的安全裕度行为，当$d_{ij}<1.5d_{\text{safety}}$时系统会记录该时刻的最小间隔比值$r=d_{\min}/d_{\text{safety}}$，作为严重程度分析的原始数据。

\textbf{冲突严重程度分级。}仅仅统计冲突次数无法反映冲突的“深浅”。为此，本文以最小间隔与安全距离的比值$r=d_{\min}/d_{\text{safety}}$为依据，将每次冲突划分为三个等级：
\begin{itemize}
    \item \textbf{轻微冲突}（$0.8\leq r<1.0$）：飞行器只是擦着安全边界经过，实际间隔仍较接近安全距离，风险可控；
    \item \textbf{中度冲突}（$0.5\leq r<0.8$）：间隔已明显侵入安全距离，需要立即采取避让动作；
    \item \textbf{严重冲突}（$r<0.5$）：间隔不足安全距离的一半，属于高风险事件，极可能演化为碰撞。
\end{itemize}

这一分级方法的价值在于：两个总冲突次数相同的模型，可能一个全是轻微擦碰、另一个频繁发生严重逼近，二者的安全性能显然不可同日而语。该分级为后续严重程度分析提供了统一口径；受本版评估周期样本量限制，按$r$分级的分布统计尚未单独报告，作为后续工作（见第6.4节）。

在此基础上，定义三个安全评估指标。无冲突完成率（CFCR）要求整个Episode中所有飞行器在零间隔违反的前提下全部通过；最小间隔违反率（MSVR）统计发生$d_{\min}<d_{\text{safety}}$的时步占全部时步的比例；平均冲突幅度（MCAE）定义为单Episode内总冲突次数除以飞行器数，即$\mathrm{MCAE}=\mathrm{TC}/N$，刻画冲突的密集程度。MCAE刻画的是冲突的“量”（平均每架飞行器的冲突次数）而非“严重程度”——严重程度信息由上述$r$分级刻画，二者不能混用。当评估规模$N$固定时，MCAE与TC严格成正比（见第5.4节的分析）。

%==========================================================================
% 4 基于DyGAT-MARL的协同管控模型
%==========================================================================
\section{基于DyGAT-MARL的协同管控模型}

\subsection{模型总体架构}

模型采用集中训练分布式执行（CTDE）范式，总体架构如图\ref{fig:architecture}所示，从下到上分四层：环境层（AirTrafficEnv，含风扰与噪声）、编码层（LSTM时序编码+ DyGAT空间聚合，含乘法门控注意力掩码）、决策层（策略头+ Q值头+GAT-Mixer，训练集中、执行分布）、训练层（课程学习+冲突损失加权+PPO）。

\begin{figure}[H]
    \centering
    \includegraphics[width=1.0\textwidth]{QQ图片20260523201329.png}
    \caption{模型总体架构——环境层、编码层（LSTM+DyGAT含乘法门控注意力掩码）、决策层（策略头+Q值头+GAT-Mixer）与训练层（课程学习+冲突损失加权+PPO），遵循集中训练分布式执行范式。}
    \label{fig:architecture}
\end{figure}

四层架构的设计思想可以概括成一句话：环境层保真，编码层求精，决策层分布，训练层集中。环境层尽量贴近真实运行；编码层通过时空注意力与物理门控，从观测里提炼决策所需的关键信息；决策层保证执行阶段的实时性与分布式特性；训练层则靠课程式学习与冲突加权，化解稀疏奖励下的优化难题。

\subsection{LSTM时序编码}

飞行器的运动有明显的时序连贯性，当前位置和速度的历史轨迹里蕴含着运动趋势，是判断未来风险的重要线索。模型用两层LSTM（隐层64维，dropout 0.1）做时序编码：
\begin{equation}
    h_i^{\text{traj}} = \text{LSTM}_{64}(o_i^{t-T+1:t}),
    \label{eq:lstm}
\end{equation}
其中$o_i^{t-T+1:t}$为最近$T=10$步的观测序列，$h_i^{\text{traj}}\in\mathbb{R}^{64}$为轨迹嵌入。LSTM的记忆单元使模型能够保留飞行器速度变化趋势与转向意图等时序信息——例如，一个正在减速避让的邻居与一个即将加速穿越的邻居，二者接下来的行为截然不同，仅凭当前瞬时状态无法区分，需要借助历史轨迹才能判断。

LSTM单元内部通过三个门控结构控制信息流：遗忘门$f_t$决定历史信息的保留比例，输入门$i_t$决定新信息的写入比例，输出门$o_t$决定隐状态的输出。其单元更新可表示为：
\begin{align}
f_t &= \sigma(W_f\cdot[h_{t-1},x_t]+b_f),\\
i_t &= \sigma(W_i\cdot[h_{t-1},x_t]+b_i),\\
\tilde{C}_t &= \tanh(W_C\cdot[h_{t-1},x_t]+b_C),\\
C_t &= f_t\odot C_{t-1} + i_t\odot\tilde{C}_t,\\
o_t &= \sigma(W_o\cdot[h_{t-1},x_t]+b_o),\quad h_t = o_t\odot\tanh(C_t),
\end{align}
其中$\sigma$为sigmoid函数，$\odot$为逐元素乘法。这套门控结构即第4.3.3节“乘法门控注意力掩码”的思想源头——区别在于，LSTM的门控信号是学习得到的，而本文的门控信号完全由物理量直接驱动。

\subsection{动态图注意力网络（DyGAT）}

DyGAT是本文的核心创新模块。它输入$N$个节点的特征（以LSTM输出$h_i^{\text{traj}}\in\mathbb{R}^{64}$为节点特征，输出维度$F_{\text{out}}=256$），输出融合了空间与时序信息、并经物理门控调制后的节点表示。模块由空间注意力、时序注意力和乘法门控注意力掩码三个组件构成。

\subsubsection{空间注意力}

空间注意力沿用GAT的标准范式，衡量邻居节点对当前节点信息的重要性。本文采用4头多头注意力机制：每个注意力头以各自的线性变换$W_h$与注意力参数$a_h$独立计算注意力分数，对各头的聚合结果取平均，以增强模型对邻居关系不同子空间的表达能力。其中$W_h\in\mathbb{R}^{F_{\text{out}}\times F_{\text{out}}}$为头$h$的特征变换矩阵，$a_h\in\mathbb{R}^{2F_{\text{out}}}$为头$h$的注意力参数。节点$i$与节点$j$在第$h$个头下的注意力分数为（其中LeakyReLU$_{0.2}$表示负半轴斜率为0.2的带泄漏修正线性单元）：
\begin{equation}
    e_{ij}^{h} = \text{LeakyReLU}_{0.2}\!\left(a_h^{\top}[W_h h_i^{\text{traj}} \parallel W_h h_j^{\text{traj}}]\right),
    \label{eq:spatial_score}
\end{equation}
其中$\parallel$为向量拼接。归一化后得到该头的注意力权重：
\begin{equation}
    \alpha_{ij}^{h} = \frac{\exp(e_{ij}^{h})}{\sum_{k\in\mathcal{N}_i}\exp(e_{ik}^{h})},
    \label{eq:spatial_softmax}
\end{equation}
节点$i$在第$h$个头下的聚合特征为$h_{i,h}^{\text{spatial}} = \sum_{j\in\mathcal{N}_i}\alpha_{ij}^{h} W_h h_j^{\text{traj}}$，最终空间特征取各头平均：
\begin{equation}
    h_i^{\text{spatial}} = \frac{1}{H}\sum_{h=1}^{H} \sum_{j\in\mathcal{N}_i}\alpha_{ij}^{h}\, W_h h_j^{\text{traj}},
    \label{eq:spatial_agg}
\end{equation}
其中$H=4$为注意力头数。空间注意力实现了由数据自适应决定“哪些邻居更值得关注”的聚合方式，但正如第2.5节所述，纯数据驱动的注意力在训练数据稀疏时可能偏离物理直觉——这为后文的乘法门控机制提供了动机。

\subsubsection{时序注意力}

为利用多步历史信息，DyGAT在注意力头上引入时序聚合。对每架飞行器，把最近$T=10$步的空间编码特征按匹配度加权求和：
\begin{equation}
    h_i^{\text{temporal}} = \sum_{\tau=1}^{T} \beta_{i,\tau}\, h_i^{\text{spatial}}(t-\tau+1),
    \label{eq:temporal_attn}
\end{equation}
其中$\beta_{i,\tau}$为时序注意力权重，由当前特征与第$\tau$个历史时刻特征的相似程度经softmax归一化得到，数值越大表示该历史时刻对当前决策的参考价值越高。时空融合表示为：
\begin{equation}
    h'_i = h_i^{\text{spatial}} + \alpha_{\text{temp}} \cdot h_i^{\text{temporal}},
    \label{eq:temporal_fusion}
\end{equation}
融合系数$\alpha_{\text{temp}}=0.3$为经验值，说明当前空间信息占主导、历史信息起辅助修正作用——这也符合“决策主要依据当前态势”的直觉。该系数作为经验超参数未做敏感性分析，这是本文超参数不确定性的来源之一，其影响留待后续工作检验。

\subsubsection{乘法门控注意力掩码（核心创新）}

上述空间与时序注意力都是纯学习机制。纯学习机制的隐患在于：当训练数据没覆盖所有交互模式时，学习注意力可能产生“反物理”的分配——比如给远距离、航向正交的飞行器很高的注意力。为从机制上约束这种风险，DyGAT用三个手工设计的指数核，把飞行器对之间的物理量编码成乘法门控信号$w_{ij}$：
\begin{equation}
    w_{ij} = \exp\!\left(-\frac{d_{ij}}{1000}\right) \cdot \exp\!\left(-\frac{|v_i - v_j|}{20}\right) \cdot \exp\!\left(-\frac{|\theta_i - \theta_j|}{\pi/4}\right),
    \label{eq:gating_w}
\end{equation}
并以乘法方式作用于各注意力头的学习注意力分数：
\begin{equation}
    e_{ij}^{h,\text{final}} = e_{ij}^{h} \odot w_{ij},\quad h=1,\ldots,H.
    \label{eq:gating_final}
\end{equation}
门控在softmax归一化之前施加，随后仍按式(\ref{eq:spatial_softmax})对邻居重新归一化，因此各节点的注意力权重之和始终为1，软性门控只改变注意力在邻居间的分配比例而不破坏归一化。由于门控信号$w_{ij}$完全由物理量驱动、与注意力头无关，同一飞行器对的所有注意力头共享同一个门控值——这保证了物理安全约束对多头注意力的全局一致性约束。

$w_{ij}$的三个因子分别对应三条基本的安全常识：距离越近，威胁越大（距离因子）；速度差越大，相对运动越剧烈，碰撞风险越高（速度差因子）；航向差越大，越接近迎面交叉，冲突可能性越大（航向差因子）。三者相乘，构成完整的门控信号。

为直观说明门控效果，表\ref{tab:gating}给出四种典型飞行器对情景下的$w_{ij}$数值（$N=21$，$\bar{v}=70$ $\mathrm{km/h}$时$d_{\text{safety}}\approx60$ m）。表中数值均由式(\ref{eq:gating_w})直接代入计算。

\begin{table}[H]
    \centering\small
    \caption{$w_{ij}$在不同飞行器对情景下的门控效果（按式(\ref{eq:gating_w})计算）}
    \label{tab:gating}
    \resizebox{0.98\textwidth}{!}{%
    \begin{tabular}{|c|c|c|c|c|c|}
        \hline
        \textbf{情景} & $d_{ij}$(m) & $|v_i-v_j|$($\mathrm{km/h}$) & $|\theta_i-\theta_j|$ & $w_{ij}$ & \textbf{门控效果} \\
        \hline
        近距离同向等速 & 100 & 5 & $\pi/12$ & $\approx0.505$ & 中度衰减（约50\%），$e_{ij}$仍可自由调节 \\
        中距离同向等速 & 300 & 10 & $\pi/8$ & $\approx0.273$ & 较强衰减（约73\%） \\
        中距离相向高速 & 200 & 40 & $\pi/2$ & $\approx0.015$ & 强衰减（98.5\%），注意力降至约1/67 \\
        远距离航向正交 & 500 & 30 & $\pi/2$ & $\approx0.018$ & 强衰减（约98\%），信息流近乎阻断 \\
        \hline
    \end{tabular}%
    }
\end{table}

表中的对比可以说明门控的区分能力：两对飞行器距离同为200 m左右时，同向等速（$\Delta v\approx0$、$\Delta\theta\approx0$）的$w_{ij}\approx0.82$，注意力几乎不受限制；相向高速（$\Delta v=40$ $\mathrm{km/h}$、$\Delta\theta=\pi/2$）的$w_{ij}\approx0.015$，注意力被压到约1/67，二者相差约55倍。而固定阈值GAT的邻接矩阵里，二者同为“邻居”，权重完全一样——这是DyGAT相对静态图方法的本质优势。表\ref{tab:gating}中的中距离相向高速与远距离航向正交两行数值接近（0.015与0.018），这是因为两行都由航向差因子（衰减到约0.135）主导，这是式(\ref{eq:gating_w})的客观计算结果，说明航向差是风险判定的主要驱动因素。

\textbf{门控方向性的说明。}直觉上，“聚焦高风险交互”容易理解为对最危险的飞行器对给予最大注意力；而式(\ref{eq:gating_w})的数学效果恰恰相反——相向高速这类碰撞几乎必然发生的飞行器对，其$w_{ij}$被压到很小（表\ref{tab:gating}中为0.015），注意力反而被抑制。这一“反直觉”设计是有意为之，理由有三。

第一，信息冗余。对相向高速的飞行器对，其冲突结局已由运动学状态几乎确定：双方相对速度持续收窄、航向正对，唯一合理的化解方式是其中一方提前减速、错峰通过。对这一对“必然相遇”的飞行器分配大量注意力，能获得的新信息有限——决策的重点不在于“要不要避让”，而在于“谁先让、让多少”。第二，软性门控保留了必要信息。$w_{ij}$只是把注意力权重缩小到约1/67而非置零，历史轨迹中关于相对速度收窄速率的信息仍然进入后续决策，模型据此选择提前减速的时机与幅度。第三，注意力重分配效应。乘法门控在softmax归一化之前施加，因此对一个高威胁对的抑制会自动把注意力“预算”让给其他仍可协商的交互对——即风险尚未确定、协同决策能真正改变结局的飞行器对。综合三点，门控的实际效果是：把有限表征能力从“结局已定、信息冗余”的交互对转移到“决策敏感、仍可协调”的交互对，同时通过保留的极小权重维持对必然冲突对的兜底感知。本文中“聚焦高风险交互对”的表述，指的是这种注意力重分配意义上的聚焦，而非对最危险对的注意力最大化。

\subsubsection{门控机制的性质分析}

乘法门控有三个重要性质，共同保证了模型的物理合理性。

\textbf{软性非阻断。}$w_{ij}\in(0,1]$，物理先验不会完全掐断信息流。即使最不利的情形（远距离+高速差+航向正交），$w_{ij}$也只是一个极小正数而非零，给学习注意力保留了在极端条件下发现非常规依赖的可能。硬性掩码“一票否决”导致的表达力损失在这里不会发生。

\textbf{非对称性。}乘法形式（$\odot$而非$+$）使门控具有非对称性：$w_{ij}$可以近乎完全压制$e_{ij}$，但$e_{ij}$永远没法“抬升”一个已经被压到接近零的权重。换句话说，物理安全约束是“学习覆盖不了”的——就算训练数据里出现反例，模型也没法通过学习推翻物理门控划定的危险边界。这一点赋予了模型安全保证的鲁棒性。

\textbf{零参数零开销。}$w_{ij}$由状态量直接计算，不需要梯度回传。在$N=21$时，单步只增加约0.2 ms的额外开销（向量化实现），相对于约213 ms的整体推理时间（第5.9节实测）几乎可以忽略，实时性基本不受影响。

从理论上看，这个门控机制本质上是\emph{结构化先验注入}（structured prior injection）。它和LSTM里的遗忘门$\sigma(Wx+b)$那种“学出来的门控”有本质区别：标准门控的门控信号是数据驱动的，本文的$w_{ij}$完全由物理量驱动、零可学习参数。更直观地说，$w_{ij}$划定了学习注意力可以操作的“有效范围”，学习注意力在这个范围内做精细化分配——二者形成“物理划界、学习微调”的互补格局。

关于三个指数核衰减常数（1000 m、20 $\mathrm{km/h}$、$\pi/4$ rad）的取值，其物理含义可以逐一解读。距离因子$\exp(-d_{ij}/1000)$的衰减常数取1000 m，大于邻域半径500 m，这意味着在可感知范围内距离因子始终保持在0.6以上，只起弱约束作用——保证远距离邻居仍保留基本的信息流，不至于被完全忽视；速度差因子$\exp(-|v_i-v_j|/20)$以20 $\mathrm{km/h}$为标尺，对应eVTOL速度范围内约七分之一的典型速度差，当速度差达到40 $\mathrm{km/h}$（两机相向飞行且均处于中等速度时很常见）该因子衰减到约0.135；航向差因子$\exp(-|\theta_i-\theta_j|/(\pi/4))$以45$^{\circ}$为标尺，航向正交（90$^{\circ}$）时衰减到约0.135。这三个常数体现了“以飞行器的典型速度尺度与航向尺度定义风险标尺”的设计思想，是手工预设的工程值，后续可通过网格搜索在更细粒度上优化，以适配不同空域尺度与机型参数。

\subsection{决策网络与GAT-Mixer混合网络}

每架飞行器维护一个单智能体网络，信息流为：LSTM编码器（2层，64维）$\to$ DyGAT层（输出256维$h_i^{\text{coupled}}$）$\to$ 策略头与Q值头。策略头是2层MLP（Tanh激活），输出动作均值$\mu_i$，动作按高斯分布$\mathcal{N}(\mu_i, 0.1^2)$采样；固定标准差0.1让策略具备可调节的探索性，又避免了方差学习的不稳定。Q值头是3层MLP（ReLU激活），输出个体Q值。策略头与Q值头共享编码层参数，好处有二：一是特征提取层参数量大幅减少，降低过拟合风险；二是策略与价值共享语义空间，有利于价值函数对策略改进的指导。

集中训练阶段，需要把各智能体信息融合成一个全局Q值$Q_{\text{tot}}$来评估联合状态-动作的质量。本文设计GAT-Mixer完成这一融合：
\begin{equation}
    Q_{\text{tot}} = \text{MLP}\!\left(\text{MeanPool}\!\left(\text{GAT}\left(\{s_i\}_{i=1}^{N}\right)\right)\right),
    \label{eq:gat_mixer}
\end{equation}
其中$\{s_i\}$为各智能体的混合网络输入向量。形式化状态$s_i\in\mathbb{R}^5$中的第5维$d_i^{\text{center}}$（到中心的距离）与位置$(x_i,y_i)$在以中心为原点的坐标系下完全冗余，因此GAT-Mixer实际使用各智能体的4维投影$[x_i,y_i,v_i,v_i^{\text{target}}]$作为输入，在全连接图上经4头多头GAT聚合（各头聚合结果取平均）后取节点均值池化，再经MLP输出。这一4维输入与3.2节5维状态定义不矛盾，差异仅在于省略了与位置冗余的$d_i^{\text{center}}$。

GAT-Mixer的设计动机有三：一是以GAT注意力替代QMIX的超网络，突破单调性约束。QMIX要求全局Q值随任一个体Q值单调不减（$\partial Q_{\text{tot}}/\partial Q_i\geq0$），这一约束简化了分解问题，但也限制了表达能力——例如，当某个智能体的个体Q值偏高、需要通过全局协调压低其权重时，单调性约束无法表达这种“个别压降”的调整。二是GAT注意力天然适应$N$的动态变化，契合课程学习中$N$从5到21的增长。三是均值池化具有平移不变性，使$Q_{\text{tot}}$对智能体排列顺序不敏感。这三条属于机制层面的设计论证，GAT-Mixer与简单均值池化混合网络、以及与QMIX单调混合网络的直接对比实验尚未单独开展，作为后续工作（见第6.4节）。此外，本文的GAT-Mixer与ICLR 2023的GraphMixer\cite{ref28}只是名字相近，后者是基于MLP-Mixer的时序链路预测模型，与本文的全局价值混合功能无关。

$Q_{\text{tot}}$与全局状态价值$V(S)$共同支撑训练目标。全局Q值以时序差分目标训练，对应损失：
\begin{equation}
    L_{q} = \mathbb{E}\left[\left(Q_{\text{tot}}(S^t,a^t) - \left(r^t + \gamma Q_{\text{tot}}(S^{t+1},a^{t+1})\right)\right)^2\right],
    \label{eq:loss_q}
\end{equation}
全局状态价值网络$V(S)$为输入各智能体状态（不足21架时补零对齐）展开向量的3层MLP，输出全局状态价值，损失为：
\begin{equation}
    L_{\text{value}} = \mathbb{E}\left[\left(V(S^t) - G^t\right)^2\right],
    \label{eq:loss_value}
\end{equation}
其中$G^t$为由GAE（第4.5.2节）计算得到的回报。

\subsection{训练机制}

\subsubsection{课程学习}

如果一开始就让21架飞行器在紧凑空域里同时学习，模型面对的奖励信号极度稀疏——绝大多数Episode中根本不发生冲突，安全奖励的梯度贡献趋近于零，模型很难从中学会如何规避危险。课程学习（Curriculum Learning）\cite{ref39}的思路即针对这一难题：先让模型在较简单的任务上掌握基础避让能力，再逐步过渡到困难任务，从而降低初始学习的难度。

本文的课程设计如下：从$N=5$架、$d_{\text{base}}=80$ m开始（宽松场景），训练过程中周期性统计最近一段时间的冲突率与完成率（本文每10个episodes统计一次），连续3个评估周期满足“冲突率低于阈值且完成率高于80\%”就触发进阶：
\begin{equation}
    N \leftarrow \min(N+2, 21), \qquad d_{\text{base}} \leftarrow \max(d_{\text{base}}-5, 50).
    \label{eq:curriculum}
\end{equation}
每次进阶增加2架飞行器、缩小5 m基本安全间隔，共9个级别。为防训练中性能回退，设有回退机制：连续5个评估周期不满足完成率要求时，$N$回退1架、$d_{\text{base}}$回升5 m。前期在$N=5$时学会的基础避让技能（减速让行、保持间隔）在后期$N=21$场景中具有迁移性，有助于缩短整体训练时间；直接的“课程学习 vs.\ 直接$N=21$训练”收敛轮次对照数据待后续补测，表\ref{tab:lambda}中不同$\lambda$的收敛周期（epoch）差异可作为训练效率的间接参考。

\subsubsection{Episode级冲突损失加权}

即使有了课程学习，冲突仍是低概率事件。本文进一步引入Episode级冲突损失加权：对包含冲突的Episode施加更大的损失权重，使模型在参数更新时给予危险样本更高的权重。

在PPO\cite{ref30} on-policy框架下，总损失为：
\begin{equation}
    L_{\text{total}} = \lambda\left(L_{\text{policy}} + L_{q} + 0.5\,L_{\text{value}}\right) + 0.01\,L_{\text{entropy}},
    \label{eq:total_loss}
\end{equation}
其中$L_{\text{policy}}$为PPO裁剪策略损失，$L_{q}$由式(\ref{eq:loss_q})给出，$L_{\text{value}}$由式(\ref{eq:loss_value})给出，$L_{\text{entropy}}$为策略熵项（权重0.01，鼓励探索、防止策略过早收敛），系数为：
\begin{equation}
    \lambda = \begin{cases}
        2.0, & \text{含冲突的Episode}\\
        1.0, & \text{无冲突的Episode}
    \end{cases}.
    \label{eq:lambda}
\end{equation}

首先说明该机制与优先经验回放（Prioritized Experience Replay，PER）的区别。PER改变的是采样分布——高优先级样本被更频繁地采样，使采样分布偏离原始分布，需要用重要性采样校正；而本文的$\lambda$只是缩放损失幅度（$\nabla_\theta(\lambda L)=\lambda\nabla_\theta L$），不涉及对数据分布的改变，因此\textbf{严格保持了PPO的on-policy特性}——所有样本仍来自当前策略，只是含冲突样本的梯度贡献被放大了2倍。

其次说明$\lambda$与PPO裁剪的交互。裁剪作用于$r_t(\theta)$层面，当某条transition的$r_t(\theta)$被clip截断时，$\lambda$对它就不再有效——被裁剪的目标值已经固定了。因此$\lambda$加权在训练早期（策略不稳定、clip截断率高）的实际效果可能弱于预期。PTR-PPO\cite{ref27}用截断重要性采样处理off-policy轨迹，是更系统的方案，可作为后续改进方向。

$\lambda$的取值直接影响训练效果，本文在$\{1.0, 1.5, 2.0, 2.5, 3.0\}$上做了实验（$n=5$），结果见表\ref{tab:lambda}。$\lambda=2.0$时TC最低（18.2）、CFCR最高（68.4\%），为最优；$\lambda=3.0$收敛最快（约300周期）但最终性能反而下降（CFCR降至57.3\%），说明权重过大导致优化振荡。当然，五个离散值的最优未必是连续空间的全局最优，更精细的调参可能带来进一步提升，但$\lambda=2.0$已足以支撑本文结论。

\begin{table}[H]
    \centering
    \caption{$\lambda$取值实验（$n=5$）}
    \label{tab:lambda}
    \begin{tabular}{|c|c|c|c|}
        \hline
        $\lambda$ & TC$\downarrow$ & CFCR(\%)$\uparrow$ & 收敛周期（epoch） \\
        \hline
        1.0 & $24.6\pm2.8$ & $53.2\pm2.3$ & $\sim800$ \\
        1.5 & $21.8\pm2.5$ & $58.7\pm2.0$ & $\sim600$ \\
        \textbf{2.0} & $\mathbf{18.2\pm2.1}$ & $\mathbf{68.4\pm1.8}$ & $\sim400$ \\
        2.5 & $19.5\pm2.3$ & $64.1\pm1.9$ & $\sim350$ \\
        3.0 & $22.3\pm3.1$ & $57.3\pm2.5$ & $\sim300$ \\
        \hline
    \end{tabular}
\end{table}

训练中的优势函数采用广义优势估计（Generalized Advantage Estimation，GAE）\cite{ref38}计算，它通过对多个时间尺度上时序差分误差的加权平均估计“在给定状态下采取某动作相对平均水平的好坏程度”，在偏差与方差之间取得平衡（本文取$\gamma=0.99$、$\lambda_{\text{GAE}}=0.95$）。综合全部设计，完整训练流程如算法\ref{alg:training}所示。

\begin{algorithm}[H]
\caption{基于DyGAT-MARL的协同管控模型训练流程}
\label{alg:training}
\begin{algorithmic}[1]
\STATE 初始化课程级别$c=1$，$N=5$，$d_{\text{base}}=80$ m，策略参数$\theta$，学习率$\eta=5\times10^{-4}$
\FOR{episode $k=1,2,\ldots,K$}
    \STATE 初始化环境（$N$架eVTOL），重置智能体隐藏状态
    \STATE $ep\_data \leftarrow \emptyset$，$has\_conflict \leftarrow \text{False}$
    \REPEAT
        \STATE 各智能体获取局部观测$o_i$，经LSTM+DyGAT编码得$h_i^{\text{coupled}}$
        \STATE 策略头输出动作均值$\mu_i$，采样动作$a_i\sim\mathcal{N}(\mu_i,0.1^2)$
        \STATE 环境执行联合动作$(a_1,\ldots,a_N)$，返回奖励$R$与冲突信息
        \IF{发生冲突} \STATE $has\_conflict \leftarrow \text{True}$ \ENDIF
        \STATE 将$(o,a,\log\pi,\hat{A},R)$存入$ep\_data$
    \UNTIL{Episode结束}
    \STATE 计算GAE\cite{ref38}优势估计（$\gamma=0.99$，$\lambda_{\text{GAE}}=0.95$）
    \STATE 计算$\lambda = 2.0 \text{ if } has\_conflict \text{ else } 1.0$
    \STATE 以$\lambda$加权总损失（式(\ref{eq:total_loss})），更新策略参数$\theta$
    \STATE 课程学习判定：连续3周期满足进阶条件则$N\leftarrow\min(N+2,21)$，$d_{\text{base}}\leftarrow\max(d_{\text{base}}-5,50)$
\ENDFOR
\STATE \RETURN 训练后的策略参数$\theta$
\end{algorithmic}
\end{algorithm}

\subsubsection{训练稳定性保障}

除了课程学习与冲突加权，训练稳定性还依赖一系列常规但必要的工程技术措施。其一是\textbf{熵正则化}：在PPO目标中加入熵项（权重0.01），鼓励策略保持适度探索，避免过早收敛到确定性的次优动作——尤其当课程学习进阶、场景突然变难时，足够的探索熵有助于策略快速适应新难度。其二是\textbf{梯度裁剪}：每步更新前对参数梯度的范数进行裁剪（上限1.0），防止个别大梯度破坏已经学好的参数。其三是\textbf{学习率指数衰减}：初始学习率$5\times10^{-4}$按指数速率衰减，使训练后期参数更新幅度逐渐减小，有利于在最优解附近精细收敛。其四是\textbf{多轮PPO内循环更新}：每个Episode收集到的数据被用于多轮（4轮）参数更新，提高数据利用效率。这些措施共同保障了第5.2节所示的三阶段平稳收敛。

\subsection{计算复杂度分析}

在低空交通管控中，算法能否实时运行直接决定了它的工程可行性，因此有必要对模型的计算复杂度进行理论分析。设单架飞行器邻域内的邻居数为$K$，节点特征维度为$F$，DyGAT各模块的计算复杂度如下：
\begin{itemize}
    \item \textbf{空间注意力}：特征线性变换为$O(KF^2)$，注意力分数计算为$O(K^2F)$；
    \item \textbf{乘法门控}：门控信号仅由物理量（距离、速度差、航向差的指数运算）直接计算，复杂度仅为$O(K^2)$，且不含任何可学习参数；
    \item \textbf{时序注意力}：复杂度为$O(TF^2)$，其中$T=10$为时序窗口长度。
\end{itemize}

由于邻域半径限制（500 m），每架飞行器的实际邻居数$K$远小于机队规模$N$（在$N=21$时平均邻居数仅约6），因此DyGAT的实际计算量随$N$近似线性增长，而非全连接图那样的$O(N^2)$增长。这一特性使模型具备良好的规模扩展性。第5.9节将给出$N=5\sim21$下实测推理时间的验证，与这里的理论分析相互印证。

%==========================================================================
% 5 实验设计与结果分析
%==========================================================================
\section{实验设计与结果分析}

本节对模型进行全面实验验证。先交代实验设置，然后从训练收敛性、综合性能、安全指标、统计显著性、消融实验、密度鲁棒性、参数敏感性与推理效率八个维度逐项分析。

\subsection{实验设置}

\subsubsection{仿真环境与参数}

实验在自研AirTrafficEnv环境中进行，全部代码基于Python 3.11与PyTorch实现，运行于NVIDIA RTX 4060（8 GB）GPU与Intel i9-14900K CPU平台。核心参数见表\ref{tab:env}。其中空域$1000\ \mathrm{m}\times1000\ \mathrm{m}$对应典型城市街区尺度，飞行器初始位置位于600 m半径的圆周上，留有充分机动空间；最大步长200步（100 s）保证即使多轮避让也有足够时间完成穿越。本文中“训练周期（epoch）”指每收集400个episodes完成一轮参数更新扫描，课程学习的进阶判定在每个评估周期（10个episodes）结束后基于最近一段时间的统计量进行。

\begin{table}[H]
    \centering\small
    \caption{仿真环境与训练超参数}
    \label{tab:env}
    \begin{tabular}{|c|c|c|c|}
        \hline
        \textbf{参数} & \textbf{取值} & \textbf{参数} & \textbf{取值} \\
        \hline
        空域范围 & $1000\,\mathrm{m}\times1000\,\mathrm{m}$ & 最大转弯角 & $\pi/6$ rad \\
        $d_{\text{base}}$（基本安全间隔） & 50 m & 时序窗口 $T$ & 10步（5 s） \\
        $\Delta t$（控制周期） & 0.5 s & 邻域半径 & 500 m \\
        最大步长 & 200步（100 s） & 训练周期（epoch） & 1000 \\
        每训练周期Episode数 & 400 & 每Episode内PPO更新轮数 & 4 \\
        目标速度 & $70\pm10$ $\mathrm{km/h}$ & $\gamma,\ \epsilon$（clip） & 0.99, 0.2 \\
        速度范围 & $[60, 140]$ $\mathrm{km/h}$ & 初始学习率$\eta$ & $5\times10^{-4}$ \\
        最大加速度 & 3.0 $\mathrm{m/s^2}$ & 最大减速度 & 4.0 $\mathrm{m/s^2}$ \\
        风扰强度 & 0.1 $\mathrm{m/s}$ & 噪声标准差$\sigma$ & 0.5 m \\
        GAE参数$\lambda_{\text{GAE}}$ & 0.95 & 熵系数 / 梯度裁剪范数 & 0.01 / 1.0 \\
        \hline
    \end{tabular}
\end{table}

为保证实验结果的统计可靠性，每个配置独立运行5个随机种子（$n=5$），每个种子在相同环境下重新初始化智能体并完成1000个训练周期；评估时冻结策略，并对所有方法使用相同的评估种子序列，确保比较条件一致。所有RL基线共享相同的训练预算与网络规模，仅在算法结构上有所区别；规则方法与工程方法直接施加其固有策略，不进行训练。这一评估流程保证了“方法间的性能差异来自算法本身、而非训练细节”这一比较前提。

仿真环境与真实低空运行的对应关系与差异。环境构建参考了：交叉路口交通流构型（对应真实UTM中的热点航路交叉）、速度自适应间隔（对应空管“速度越高间隔越大”的原则及DAA well clear概念）、风扰与传感器噪声（对应真实气象与定位误差）。与公开基准相比，BlueSky\cite{ref52}提供开源空管仿真框架、OpenSky\cite{ref53}提供真实ADS-B轨迹数据、SUMO\cite{ref54}提供微观交通流仿真，均可作为环境有效性对标或参数标定的来源。由于本文聚焦算法层面的机制验证，暂未与上述基准逐场景对标，该对标留待后续工作；环境参数（空域尺度、速度范围、间隔、控制周期）均取自低空eVTOL运行的常见量级，具体来源见第3章与表\ref{tab:env}。

\subsubsection{基线方法}

选取9种覆盖“规则调度—单/独立RL—CTDE-MARL—工程安全方法”四个层级的基线。规则调度包括轮询（RR）和先到先服务（FCFS），代表无学习能力的最朴素策略，用来衡量学习方法的增益下界；单/独立RL包括集中式DQN和独立A2C；CTDE-MARL包括MADDPG、COMA、QMIX和MAPPO，其中MAPPO是最强基线；工程方法为ORCA\cite{ref4}。ORCA的前身VO在$N=21$下被ORCA一致主导，主实验不再单独报告。

所有RL基线与本文方法共享相同的环境与奖励函数，保证比较公平；规则与工程方法直接施加其固有策略。需要交代离散动作基线的适配方式：DQN、COMA与QMIX的输出为离散动作，本文将连续加速度空间$a_i\in[-4,3]$ $\mathrm{m/s^2}$按1 $\mathrm{m/s^2}$步长离散化为8档动作（$-4,-3,\ldots,3$），其余观测、奖励与环境均与连续方法完全一致；独立A2C与MADDPG直接输出连续加速度。离散化粒度对应的单步速度变化约1.8 $\mathrm{km/h}$，与连续方法在动作分辨率上的差异是影响对比公平性的因素之一，已在本节讨论中说明。

工程方法ORCA的适配细节如下：采用速度障碍互惠避碰框架，安全半径取当前时刻的速度自适应安全距离$d_{\text{safety}}$，即每个控制步按当前平均速度重算避碰约束，以适配本文的自适应间隔设定；首选速度指向交叉中心、限幅于$[60,140]$ $\mathrm{km/h}$。由于ORCA假设圆形智能体且本身不含最大加/减速与航向约束，本文将ORCA输出的速度增量进一步限幅到与学习智能体相同的加速度区间$[-4,3]$ $\mathrm{m/s^2}$，并使其遵循相同的航向转向约束，以保证与学习方法的动力学可比性。自适应安全半径是对ORCA固定半径假设的一种近似，该近似的影响留待后续用固定半径对照实验检验。

MPC和CBF未纳入比较，原因有二：MPC在21架飞行器场景下实时性难以满足0.5 s控制周期；CBF在高密度下安全集合交集容易退化。二者与MARL的混合架构是重要的未来方向。

\subsubsection{评价指标}

指标分“表层”和“深层”两个层级。表层指标刻画整体运行性能，包括累计奖励（CR，整个Episode的折扣累计奖励）、总冲突次数（TC，整个Episode所有时刻所有飞行器对之间的冲突数量累计）、平均延误（AD，相对理想飞行时间的平均延迟）。深层指标刻画安全性能：通过率（PCR，成功通过交叉中心的飞行器占比）、无冲突完成率（CFCR，全程零冲突且全部通过的Episode占比）、最小间隔违反率（MSVR，发生$d_{\min}<d_{\text{safety}}$的时步占比）、平均冲突幅度（MCAE，单Episode内总冲突次数除以飞行器数，即$\mathrm{TC}/N$，刻画冲突密度）。

这里需要区分PCR与CFCR：PCR只衡量飞行器是否到达，CFCR则要求飞行器在不发生任何间隔违反的前提下到达。若只报告PCR，可能系统性高估模型的安全性能——这是本文引入CFCR的核心动机。

\subsection{训练收敛性分析}

训练过程中的累计奖励、总冲突次数与课程进阶情况如图\ref{fig:training}所示，整体呈明显的三阶段特征。

\begin{figure}[H]
    \centering
    \includegraphics[width=0.8\textwidth]{QQ图片20260504231524.jpg}
    \caption{训练收敛曲线。阴影表示$\pm1$标准差（$n=5$），红色虚线标记课程学习进阶点。}
    \label{fig:training}
\end{figure}

\textbf{阶段I（0--200周期，快速上升期）：}累计奖励从约$-2000$升到$+800$，总冲突次数从80降到25。此时课程学习处于初级阶段（$N=5\sim11$），场景宽松，模型快速掌握基本避让策略。

\textbf{阶段II（200--600周期，进阶波动期）：}课程学习触发$N:11\to19$的多次进阶。每次进阶后，新加入的飞行器让场景密度骤增，模型出现1--3个周期的短暂回落（图中红色虚线标记的进阶点），随后迅速适应。这种“进阶—回落—再适应”的锯齿形曲线，是课程学习正常工作的标志。

\textbf{阶段III（600--1000周期，收敛稳定期）：}累计奖励稳定在$1286.7\pm18.2$，总冲突次数收敛到$18.2\pm2.1$，在$N=21$的高密度场景下达到稳定性能。

收敛曲线也是训练稳定性的直接体现：若曲线剧烈振荡或长期不收敛，通常意味着训练机制存在问题。本文的三阶段收敛特征说明课程学习有效缓解了稀疏奖励导致的训练不稳定，为最终性能奠定了基础。

进一步观察三阶段中奖励构成的变化可以发现：阶段I的奖励上升主要来自安全奖励的改善——模型在$N=5$的宽松场景下迅速学会了拉开安全间隔；阶段II的波动则主要来自课程进阶瞬间冲突率的短暂回升；进入阶段III后，奖励的稳定意味着安全、速度、协同与效率四部分奖励之间达成了平衡。这一平衡的达成并非偶然：协同奖励抑制了个体的争先行为，效率奖励保证了通行效率，二者与安全奖励共同作用，才使得模型在高密度下既能保证安全、又不至于过度保守。

\subsection{综合性能对比}

表\ref{tab:comparison}给出21架eVTOL场景下各方法的综合性能（$n=5$，均值$\pm$标准差）。

\begin{table}[H]
    \centering
    \caption{综合性能对比（21架eVTOL，$n=5$）}
    \label{tab:comparison}
    \begin{tabular}{|c|c|c|c|c|}
        \hline
        \textbf{方法} & \textbf{CR$\uparrow$} & \textbf{TC$\downarrow$} & \textbf{AD(s)$\downarrow$} & \textbf{PCR(\%)$\uparrow$} \\
        \hline
        RR & $-1256.3\pm89.2$ & $42.6\pm5.3$ & $1.58\pm0.12$ & $89.2\pm3.5$ \\
        FCFS & $-1187.5\pm76.4$ & $38.9\pm4.7$ & $1.42\pm0.10$ & $91.5\pm2.8$ \\
        DQN & $-892.5\pm65.3$ & $28.3\pm3.9$ & $1.21\pm0.09$ & $94.7\pm2.1$ \\
        独立A2C & $-657.8\pm52.6$ & $19.7\pm2.8$ & $0.93\pm0.07$ & $96.3\pm1.7$ \\
        ORCA & $-198.7\pm35.8$ & $22.8\pm2.9$ & $0.76\pm0.06$ & $96.0\pm1.6$ \\
        MADDPG & $-215.6\pm41.2$ & $25.4\pm3.2$ & $0.87\pm0.06$ & $95.8\pm1.9$ \\
        COMA & $187.3\pm35.7$ & $22.1\pm2.9$ & $0.79\pm0.06$ & $97.1\pm1.5$ \\
        QMIX & $426.5\pm28.9$ & $23.1\pm2.7$ & $0.76\pm0.05$ & $97.5\pm1.3$ \\
        MAPPO & $892.4\pm23.5$ & $20.3\pm2.4$ & $0.65\pm0.04$ & $98.2\pm1.1$ \\
        \textbf{本文} & $\mathbf{1286.7\pm18.2}$ & $\mathbf{18.2\pm2.1}$ & $\mathbf{0.57\pm0.03}$ & $\mathbf{99.1\pm0.8}$ \\
        \hline
    \end{tabular}
\end{table}

从表中可以归纳出四条结论。第一，学习方法对规则方法呈系统性优势：RR的TC高达42.6、AD超过1.4 s、奖励为负，基本处于勉强运行的状态；本文相对RR将TC降低57.2\%、AD缩小63.9\%。第二，CTDE-MARL全面优于单智能体方法（DQN、独立A2C），再次验证集中训练对协作的重要性。第三，本文方法在四项指标上全部最优。第四，ORCA（TC=22.8）明显强于规则方法，但它的CFCR只有约50\%（见表\ref{tab:safety}），低于本文的68.4\%——ORCA缺乏多步轨迹预测能力，高密度下速度空间的可行解容易退化。

从方法族的角度进一步看，四类基线的表现差异本身也很有启发性。规则方法（RR、FCFS）表现最差，是因为它们把“多机协同”退化成了“单机调度”——轮询或按到达顺序放行，完全没有考虑飞行器间的相对运动状态，一旦密度升高，积压的飞行器会同时涌入路口，冲突不可避免。单智能体方法（DQN、独立A2C）虽然学会了避让行为，但每个智能体都在与环境“单方面”互动，无法感知其他智能体的意图，学习到的策略在博弈耦合下趋于保守或振荡，因此性能优于规则方法、劣于CTDE方法。CTDE方法（MADDPG、COMA、QMIX、MAPPO）通过集中训练显式建模了智能体间的协作关系，性能整体上了一个台阶，其中以PPO为基础的MAPPO凭借稳定的策略优化过程表现最佳。工程方法ORCA在中等密度下表现不俗，但缺乏对整体交通流的多步预测，高密度下速度空间可行解的退化，最终体现在其CFCR只有约50\%上。这四条线索共同刻画了本文方法所处的性能位置：它在CTDE范式的基础上，通过物理门控进一步补上了纯学习范式在时空耦合建模上的短板。

关于本文与ORCA的TC差异（18.2 vs.\ 22.8，降幅20.2\%）需要谨慎解读：以$n=5$和标准差估算，这个差异的统计显著性可能处于边界状态，系统性的比较要等$n\geq20$的样本。

\subsection{安全指标专项分析}

表层指标与深层安全指标的对比见表\ref{tab:safety}。

\begin{table}[H]
    \centering
    \caption{安全指标对比（$n=5$）}
    \label{tab:safety}
    \begin{tabular}{|c|c|c|c|c|}
        \hline
        \textbf{方法} & \textbf{CFCR(\%)$\uparrow$} & \textbf{MCAE$\downarrow$} & \textbf{MSVR(\%)$\downarrow$} & \textbf{PCR(\%)$\uparrow$} \\
        \hline
        RR & $12.3\pm4.1$ & $2.03\pm0.25$ & $8.7\pm1.2$ & $89.2\pm3.5$ \\
        FCFS & $18.7\pm3.6$ & $1.85\pm0.22$ & $7.4\pm1.0$ & $91.5\pm2.8$ \\
        ORCA & $50.2\pm3.0$ & $1.09\pm0.14$ & $3.8\pm0.6$ & $96.0\pm1.6$ \\
        DQN & $37.5\pm3.2$ & $1.35\pm0.19$ & $5.1\pm0.8$ & $94.7\pm2.1$ \\
        独立A2C & $48.2\pm2.8$ & $0.94\pm0.13$ & $3.6\pm0.6$ & $96.3\pm1.7$ \\
        MADDPG & $42.6\pm3.0$ & $1.21\pm0.15$ & $4.2\pm0.7$ & $95.8\pm1.9$ \\
        COMA & $51.3\pm2.7$ & $1.05\pm0.14$ & $3.8\pm0.6$ & $97.1\pm1.5$ \\
        QMIX & $53.7\pm2.5$ & $1.10\pm0.13$ & $3.5\pm0.5$ & $97.5\pm1.3$ \\
        MAPPO & $57.8\pm2.2$ & $0.97\pm0.11$ & $3.0\pm0.4$ & $98.2\pm1.1$ \\
        \textbf{本文} & $\mathbf{68.4\pm1.8}$ & $\mathbf{0.87\pm0.10}$ & $\mathbf{2.4\pm0.3}$ & $\mathbf{99.1\pm0.8}$ \\
        \hline
    \end{tabular}
\end{table}

\textbf{PCR与CFCR之间约30个百分点的系统性差距，是本文最重要的实证发现。}以本文方法为例，PCR高达99.1\%——几乎全部飞行器都完成了穿越——但CFCR只有68.4\%，意味着约三成Episode存在至少一次间隔违反。这个现象在MAPPO上更明显（PCR 98.2\% vs.\ CFCR 57.8\%，差约40个百分点），说明它不是本文方法的特例，而是CTDE-MARL在安全攸关评估中的共性挑战。PCR与CFCR是两种口径不同的指标：PCR在飞行器层统计成功通过的比例，CFCR在Episode层要求全程零违反且全部通过。二者分母不同，直接相减得到的“约30个百分点差距”用于突出CFCR的严格性，其量级主要源于CFCR“全程所有飞行器零违反”的极严格定义，而非飞行器层面的“未到达”。

这个发现的方法论意义在于：如果安全评估只报PCR，会系统性高估模型的安全性能。CFCR把“只要不出事故”的宽松标准，提升为“全程零风险接触”的严格标准，更贴合低空交通安全的实际诉求。

本文方法在MCAE指标上为0.87，低于MAPPO的0.97。本文的MCAE定义为总冲突次数除以飞行器数（$\mathrm{MCAE}=\mathrm{TC}/N$），在评估规模$N=21$固定的条件下，表\ref{tab:safety}中所有方法的MCAE值都等于其TC除以21（例如本文$18.2/21=0.87$、MAPPO$20.3/21=0.97$），因此MCAE与TC严格成正比，刻画的是冲突密度而非冲突严重程度。基于第3.5节$r$分级的各严重等级（轻微/中度/严重）分布统计，受本版样本量限制未单独报告，作为后续工作。从机制层面看，门控对“相向高速”等高威胁交互对的注意力进行软性抑制，使模型倾向于提前、平滑地化解危险遭遇；这一行为是否同时降低了已发生冲突的严重程度（而非仅减少冲突次数），需要上述严重度分布分析加以验证，是后续实验的重点。

\subsection{统计显著性检验}

为确认本文方法与最强基线MAPPO的差异不是随机波动，用Welch $t$检验\cite{ref42}（双侧，显著性水平$\alpha=0.05$，$n_1=n_2=5$）做统计检验。Welch $t$检验适用于两组样本方差可能不等的独立样本均值比较。除$p$值外，还报告Cohen's $d$效应量\cite{ref43}与统计功效。效应量以标准差为单位刻画两组均值差距的大小，是量纲无关的指标，一般$d>0.8$即视为大效应；统计功效指当真实差异确实存在时，检验能够正确发现该差异的概率，功效越高，越不容易漏报真实效应。结果见表\ref{tab:ttest}。

\begin{table}[H]
    \centering
    \caption{本文 vs.\ MAPPO统计检验与功效分析}
    \label{tab:ttest}
    \begin{tabular}{|c|c|c|c|c|c|c|}
        \hline
        \textbf{指标} & \textbf{均值差} & \textbf{$t$(df=8)} & \textbf{$p$} & \textbf{$d$} & \textbf{Power} & \textbf{结论} \\
        \hline
        TC & $-2.1$ & 1.47 & 0.180 & 0.93 & 38\% & 不显著（$n\geq20$需） \\
        CFCR & $+10.6\%$ & 8.34 & $<0.001$ & 5.27 & $>99\%$ & \textbf{高度显著} \\
        MCAE & $-0.10$ & 1.50 & 0.172 & 0.95 & 39\% & 不显著（$n\geq20$需） \\
        MSVR & $-0.6\%$ & 2.68 & 0.029 & 1.69 & 70\% & \textbf{显著} \\
        \hline
    \end{tabular}
\end{table}

关于统计检验的前提，需要作三点说明。其一，所有方法共享同一评估种子序列，因此本文方法与MAPPO是在同一批随机场景上评估的，严格意义上构成配对设计。本文仍采用独立样本的Welch $t$检验，理由如下：共享场景可能使两种方法的指标产生正相关，此时配对检验的标准误小于独立检验，独立检验给出的标准误偏大、$p$值偏大，因此本文报告的显著性结论是保守口径，不会被高估；且两种方法的策略网络独立初始化、独立优化，指标中还包含各自训练的随机波动，并非严格配对。作为稳健性检验，待原始per-seed数据公开后将以配对置换检验复核，预期对CFCR/MSVR的显著结论只会增强。其二，CFCR与MSVR为$[0,1]$区间（百分比）的比例型指标，其数值来自每个seed在数百个评估episode上的平均，由中心极限定理近似服从正态分布，但在$n=5$下正态性难以严格保证。对此，本文同时报告效应量与功效作为补充，且CFCR的$t=8.34$、$p<0.001$远在阈值之外，轻微的非正态性不会改变显著结论；以bootstrap或置换检验复核比例指标的正态性假设列入后续工作。其三，MCAE与TC的信息高度重叠（见第5.4节），二者的检验结论也相互印证，均受$n=5$限制未达显著。

\textbf{核心统计结论：}本文在CFCR（$p<0.001$，效力$>99\%$）和MSVR（$p=0.029$，效力70\%）两项安全指标上取得了统计可信的改善，这是本文的核心统计贡献。CFCR的效应量$d=5.27$远超“大效应”阈值0.8，说明$+10.6\%$的提升不仅在统计上显著，工程意义上也相当可观。

TC和MCAE的差异（$d=0.93$与0.95）是大效应量，但在$n=5$下统计效力只有约38\%，未达显著（$p=0.180$与0.172）。这里的$p=0.180$要恰当解读为“在现有样本量下没检测到显著差异”，而不是“确认无差异”——大效应量恰恰暗示真实效应很可能存在。据功效分析，需$n\geq20$才让检验达到80\%效力。这个区分对科学解读至关重要：既不能把不显著夸大成“确认无差异”，也不能忽视它的方向性证据价值。

效应量的实际意义不应当只由经验阈值来判断。以CFCR为例，$d=5.27$配合$t=8.34$说明：在5个种子的重复实验中，本文与MAPPO在无冲突完成率上的差异始终稳定地朝同一方向出现，几乎不存在偶然波动的空间。而TC虽然效应量也达到0.93，但两组样本的分布存在较大重叠，个别种子中MAPPO可能表现更好——这即$p=0.180$所反映的“证据不足”，而非“确定无差异”。区分“效应量（差异有多大）”与“证据强度（证据有多足）”这两个维度，有助于更准确地阅读本文的结果，也提示了增大样本量的实际必要性。

\subsection{消融实验}

为厘清各模块的独立贡献，将完整模型与移除单一模块的变体对比，结果见表\ref{tab:ablation}。

\begin{table}[H]
    \centering
    \caption{消融实验（$n=5$）}
    \label{tab:ablation}
    \resizebox{0.92\textwidth}{!}{%
    \begin{tabular}{|c|c|c|c|c|}
        \hline
        \textbf{配置} & \textbf{CR$\uparrow$} & \textbf{TC$\downarrow$} & \textbf{AD(s)$\downarrow$} & \textbf{PCR(\%)$\uparrow$} \\
        \hline
        完整模型 & $\mathbf{1286.7\pm18.2}$ & $\mathbf{18.2\pm2.1}$ & $\mathbf{0.57\pm0.03}$ & $\mathbf{99.1\pm0.8}$ \\
        \hline
        移除DyGAT（整体替换为无门控的标准GAT） & $215.3\pm32.6$ & $49.8\pm4.5$ & $0.79\pm0.05$ & $94.3\pm1.9$ \\
        \hline
        移除课程学习（直接$N=21$） & $587.2\pm27.4$ & $28.4\pm3.1$ & $0.68\pm0.04$ & $97.2\pm1.4$ \\
        \hline
        移除冲突损失加权（$\lambda\equiv1.0$） & $763.5\pm24.1$ & $24.6\pm2.8$ & $0.65\pm0.04$ & $97.8\pm1.2$ \\
        \hline
        同时移除课程学习与冲突加权 & $128.6\pm38.5$ & $35.7\pm3.9$ & $0.82\pm0.06$ & $95.6\pm1.7$ \\
        \hline
    \end{tabular}%
    }
\end{table}

消融结果传递了三条信息。第一，\textbf{DyGAT模块整体贡献最大}：将其整体移除（退化为无门控的标准GAT，同时去掉空间门控、时序融合与动态边权重）后，TC从18.2上升到49.8（+173.6\%），CR从1286.7骤降到215.3，直接验证了DyGAT模块的核心价值。DyGAT由空间注意力、时序注意力与乘法门控三个组件构成，+173.6\%反映的是整个模块的贡献，其中乘法门控的独立贡献尚需“仅$w_{ij}$/仅$e_{ij}$/$e_{ij}\odot w_{ij}$”的组件级消融加以分解；受训练预算限制，组件级消融未在本版完成，作为最高优先级的后续实验（见第6.4节）。第二，课程学习与冲突加权的独立贡献：单独去掉课程学习，TC增加56.0\%；单独去掉冲突加权，TC增加35.2\%。第三，二者的联合效应：同时去掉时TC增加96.2\%，接近二者独立效应的加和（56.0\%+35.2\%=91.2\%）。这一结果只能说明两机制未出现明显的正/负交互，不能严格推出“作用独立”——二者叠加后可能共享部分安全收益，且$n=5$下96.2\%与91.2\%之间约5个百分点的差异没有置信区间、无法作显著性判断。因此更稳妥的表述是：课程学习与冲突加权的联合效应与其独立效应之和相当，未观察到明显交互，可作为两者可叠加的经验证据。

\subsection{密度鲁棒性分析}

前文实验都在$N=21$下进行。为检验密度鲁棒性，把飞行器数量从5逐步加到21，观察各方法TC的变化趋势（图\ref{fig:density}）。三条线的增长斜率（本文约0.75次/架，MAPPO约1.15次/架，RR约2.4次/架）由各密度点上的TC均值线性拟合得到。该拟合基于4个密度点且$n=5$，样本量不足以给出斜率的显著性检验与置信区间，此处仅作为趋势描述而非统计推断。

从$N=5$、10、15到21的递进过程，可以更细致地观察这种差距的累积效应。在$N=5$的低密度阶段，空域相对宽裕，所有方法都能维持较低的冲突水平，方法间差距很小；当$N$增大到10，飞行器对的交互数量从10对增至45对，冲突开始频繁出现，学习方法的优势逐渐拉开；当$N$达到15以上，全连接图上的注意力信息急剧膨胀，缺乏门控机制的方法开始被大量弱相关邻居稀释有效信息，冲突增速明显加快；到$N=21$时，交互对达到210对，三条线的斜率差异被完全放大，本文方法的优势也最为显著。这一由密度累积效应驱动、逐级放大的性能差距，说明本文方法的优势并非来自某个特定密度下的偶然结果，而是随着场景难度上升持续扩大的系统性优势。

\begin{figure}[H]
    \centering
    \includegraphics[width=0.65\textwidth]{QQ图片20260524171218.png}
    \caption{不同密度下冲突次数对比。误差条表示$\pm1$标准差（$n=5$）。}
    \label{fig:density}
\end{figure}

斜率对比的工程含义很直观：密度每增加一架，本文只多约0.75次冲突，而规则方法要多2.4次。低密度时差距还不明显，一旦逼近真实运行场景的极限密度，差距会线性放大。DyGAT在高密度下的优势可归因于门控机制：邻居数量增多后，全连接图上的注意力信息急剧膨胀，无门控的方法容易被大量弱相关邻居稀释有效信息，而乘法门控将注意力集中于决策敏感的少数飞行器对，使模型在高密度下仍能保持对关键交互的聚焦。

\subsection{参数敏感性分析}

测试了三个关键超参数：基本安全间隔$d_{\text{base}}$、风扰强度与传感器噪声（图\ref{fig:sensitivity}）。

\begin{figure}[H]
    \centering
    \includegraphics[width=0.8\textwidth]{QQ图片20260524162300.png}
    \caption{超参数敏感性分析。红色虚线标记默认参数取值。误差条表示$\pm1$标准差。}
    \label{fig:sensitivity}
\end{figure}

\textbf{安全间隔-效率权衡：}$d_{\text{base}}$增大使TC降34.7\%（更安全），但AD升21.2\%（更延误）。这符合物理直觉：更大的间隔意味着更大的缓冲，但也限制了交通流的紧凑性。默认50 m在二者之间取了平衡。

\textbf{风扰的破坏性：}风扰强度增大使TC上升约28\%。风扰破坏航迹精确性，迫使模型频繁修正速度航向，冲突概率上升。

\textbf{噪声的有限影响：}传感器噪声增大使TC上升约15\%，三者中最小。原因在于LSTM对时序信息的平滑作用——历史序列的聚合天然过滤单步噪声，体现了时序建模的鲁棒性价值。

从物理机理上解释三项敏感性结果也颇有意义。安全间隔的权衡，本质上反映了“安全”与“容量”这一对空管领域的永恒矛盾：间隔越大，单架飞行器占用的空域越多，可容纳的交通流密度越低，这在$d_{\text{base}}$增大时表现为TC下降与AD上升的双向变化。风扰的影响则体现了动力学不确定性对协同策略的破坏：风扰使飞行器的实际位置偏离指令轨迹，模型被迫把一部分注意力从“协同”转向“纠偏”，冲突概率随之上升。噪声影响最小这一现象，可以从LSTM的时序平滑特性得到解释——历史窗口内多个时刻的观测被联合编码，单步随机噪声在聚合过程中被平均化，体现了时序建模天然的抗噪能力。敏感性分析的意义在于确认性能的稳健边界：参数发生$\pm20\%$扰动时性能不会剧烈退化，说明性能并非依赖过拟合特定参数组合。本节的敏感性分析覆盖了$d_{\text{base}}$、风扰与噪声三类参数，未覆盖最大加速度$a_{\max}$与时序融合系数$\alpha_{\text{temp}}$，这两类参数的敏感性留待后续实验。

\subsection{推理效率分析}

低空管控对实时性要求苛刻：控制周期0.5 s，单步决策时间必须远小于它。各方法在不同$N$下的单步推理时间见图\ref{fig:computation}。

\begin{figure}[H]
    \centering
    \includegraphics[width=0.8\textwidth]{QQ图片20260524163903.jpg}
    \caption{不同飞行器数量$N$下单步推理时间对比（RTX 4060 GPU）。}
    \label{fig:computation}
\end{figure}

本文方法在$N=21$时单步推理实测约213 ms（RTX 4060 GPU，含LSTM/DyGAT前向，采用向量化的门控计算），相对0.5 s的控制周期有约2.3倍的实时性余量；推理时间随$N$近似线性增长（$N=5/10/15/21$实测约15/46/96/213 ms/步）。需要说明，上述实测来自当前逐智能体串行前向的实现，其延迟主要来自各智能体LSTM编码与DyGAT聚合中的重复GPU kernel启动开销；将全部智能体的前向合并为一次批量化的GPU调用，可将单步延迟压缩一个量级（降至几十毫秒量级），此时即便计入UTM部署中50--200 ms的通信与感知时延，整个控制周期内仍有充足的计算余量。因此，本文方法的实时性瓶颈不在于模型本身，而在于当前实现未做前向批量化优化，这是工程部署中可低成本解决的环节。乘法门控的计算开销极低：它只涉及物理量的指数运算、无需梯度回传，在$N=21$时单步仅增加约0.2 ms（向量化实现，相对213 ms整体推理可忽略）——即门控机制以可忽略的额外推理开销换取可度量的安全性能提升。

\subsection{典型穿越情景的定性分析}

为进一步展示模型在典型穿越情景中的行为，本节以两种常见的遭遇构型作简要的定性说明。

\textbf{构型一：相向高速遭遇。}两架eVTOL从路口两侧相向驶来，相对速度大、冲突风险极高。按门控公式，这类飞行器对的$w_{ij}$约0.015，学习注意力被强烈衰减；该构型的化解机制在第6.2节有完整论述。

\textbf{构型二：同向追近遭遇。}后方飞行器速度高于前方，若持续保持则可能在路口附近追近。此时两机速度差是主要风险来源，门控的速度差因子起主导抑制作用，模型倾向于让后方飞行器收敛到与前机相近的速度，保持稳定间距后随队通过。

上述为基于门控信号与训练奖励的机制推演，对单次运行轨迹的逐帧回放分析留待后续工作。

%==========================================================================
% 6 讨论与结论
%==========================================================================
\section{讨论与结论}

\subsection{与MS-GRL的比较}

Li等\cite{ref10}提出的MS-GRL是图增强MARL在低空无人机冲突消解领域的代表性工作，与本文具有较强的可比性，三个层面的差异值得逐一说明。算法上，MS-GRL用Double DQN（离散动作空间），本文用PPO（连续加速度）；连续动作能输出更平滑的加速度指令，避免离散动作导致的航迹抖动。图结构上，MS-GRL用固定距离阈值建图，凡在阈值半径内一律当邻居；本文用乘法门控的动态边权重，边权重随物理状态连续变化。场景规模上，MS-GRL处理100架UAV的编队冲突，本文处理21架eVTOL的交叉路口，前者是“大规模稀疏冲突”，后者是“小规模高密度冲突”，挑战重心不同。

门控的区分能力可以落到具体数字上：两对距离均为200 m的飞行器，一对同向等速（$\Delta v\approx0$、$\Delta\theta\approx0$）时$w_{ij}\approx0.82$，一对相向高速（$\Delta v=40$ $\mathrm{km/h}$、$\Delta\theta=\pi/2$）时$w_{ij}\approx0.015$，注意力权重相差约55倍（$0.82/0.015\approx55$），同时相向高速对的注意力相对其无门控上限（$w=1$）降至约1/67。固定阈值GAT把两者一视同仁，完全捕捉不到这个差异；DyGAT却实现了对“同样距离、不同风险”飞行器对的自适应区分。MS-GRL与本文在路线上互补，把前者的多尺度图框架与本文的物理门控思想结合，是未来值得尝试的方向。

一个值得讨论的问题是，本文方法在更大规模场景（如$N=100$架）下的可扩展性。从计算角度看，DyGAT的邻域限制（500 m）保证了每架飞行器的邻居数$K$不随$N$线性增长，因而单步计算量随$N$近似线性扩展（第4.6节的理论分析已给出依据）；从方法角度看，MS-GRL已验证了图增强方法在$N=100$规模下的可行性。因此，将本文的物理门控思想嵌入多尺度图框架、以承接更大规模的机队，是迈向百架级低空交通管控的自然路径。

\subsection{核心发现与机制解读}

\textbf{CFCR揭示的安全评估盲区。}PCR（99.1\%）与CFCR（68.4\%）之间约30个百分点的差距，在MAPPO上更大（约40个百分点），说明只报PCR会系统性高估CTDE-MARL的安全性能。这个发现不限于本文方法，而是指出了该领域评估方法论的一个普遍盲区。

\textbf{绝对安全水平与工程可接受性。}本文以“优于基线”作为相对论证，但作为UTM组件，模型达到何种绝对安全水平才算可接受，需要与民航安全目标对标。本文CFCR=68.4\%，意味着约1/3的Episode仍发生至少一次间隔违反；但“间隔违反”与“碰撞”有本质区别（第2.1节），在低空eVTOL场景下，间隔违反的风险严重程度由第3.5节的$r$分级刻画，绝大多数违反属轻微或中度（该分布统计待后续补充）。从工程对标看，ICAO Doc 4444\cite{ref49}给出了传统空管的安全目标框架，RTCA DO-365\cite{ref51}与JARUS SORA\cite{ref50}分别从探测避让能力与运行风险评估角度给出了无人机系统的安全要求；把CFCR这类Episode级零违反指标与上述安全目标框架建立明确的量化对应（例如将“严重冲突率”对标到某一可接受的事故率水平），是本文方法走向工程可接受性论证的必要一步，属于后续工作。本文在此明确表态：相对优势论证与绝对安全对标是两件不同的事，前者已由统计检验支撑，后者需要更系统的标准对接。

\textbf{乘法门控的贡献结构。}DyGAT整体替换让TC增加173.6\%，但其内部三个组件的独立贡献尚未分解。尤其是第4.3.3节提出的假设——$w_{ij}$在学习注意力产生偏差时提供必要约束——引出一个待验证的问题：如果只用$w_{ij}$做注意力（完全不学$e_{ij}$），性能会接近完整DyGAT吗？回答它需要“仅$w_{ij}$ / 仅$e_{ij}$ / $e_{ij}\odot w_{ij}$”的三路组件级消融。

\textbf{统计证据的差异化图景。}本文的核心统计贡献定位于CFCR（$p<0.001$，$d=5.27$）和MSVR（$p=0.029$，$d=1.69$）的显著改善，这两项在安全攸关评估里工程意义更直接。TC的差异（$p=0.180$，$d=0.93$）应视为方向性证据——大效应量说明真实效应很可能存在，但$n=5$的统计效力不足以确认。把论述重心从TC移到CFCR/MSVR，是对统计证据强度的诚实反映。

\textbf{机制层面的直观理解。}为从机制上理解本文方法为何有效，可以借助一个典型的“相向高速”遭遇场景作定性说明。按门控公式，这类飞行器对的$w_{ij}$约0.015，学习注意力被强烈衰减——模型不会在这对“必然相遇”的飞行器上浪费注意力做盲目试探，而是依据历史轨迹信息识别出双方相对速度正在持续收窄，从而由其中一方提前小幅减速、另一方保持航向，以错峰方式依次通过路口。这种“提前、平滑、小幅”的化解方式，与ORCA在互惠假设下双方同时大幅减速、效率损失明显的表现形成对比。模型并非被显式“教会”这一行为，而是在门控机制与协同奖励的共同作用下自然涌现的结果：门控把注意力集中到决策敏感的交互对上，协同奖励则惩罚与邻居的速度失配，二者的耦合使“主动让行”成为性价比最高的策略。

\subsection{与工程安全方法的衔接讨论}

本文的统计证据集中在CFCR与MSVR两项安全指标上，MPC与CBF作为形式化安全保证方法的代表，尚未纳入实验对比。学习方法与工程方法并非相互对立：MARL在复杂高密度场景中的适应性优势，与CBF在安全边界上的严格形式化保证，天然互补。一个合理的混合架构是让MARL负责常态下的高效协同决策，由CBF作为安全监督层，在策略输出可能越界时进行兜底修正——这种“学习决策+安全兜底”的运行时保障（runtime assurance）思想，是低空交通管控走向工程落地的重要方向。同样值得期待的是MARL与MPC的融合：以MPC的预测能力为MARL提供约束化的动作修正，或反过来以MARL的自主学习弥补MPC对复杂交互建模的不足。本文的消融与统计结果已经表明，在纯学习框架内部，DyGAT所注入的物理先验就能显著提升安全性，这为“在学习框架内部融入工程安全思想”提供了经验支持——物理先验既可以放在门控信号里，也可以放在安全监督层中，两种形式并不冲突。

\subsection{结论}

本文针对城市低空十字交叉路口的eVTOL协同管控问题，提出了一种融合乘法门控注意力掩码的耦合多智能体强化学习模型。核心创新是DyGAT中基于物理量的乘法门控机制——用距离、速度差和航向差编码出软性注意力掩码，实现飞行器间时空耦合关系的细粒度建模。配套的速度自适应安全距离、分层奖励、课程学习与冲突损失加权，共同解决了稀疏冲突样本下的训练困难。

在21架eVTOL的二维交叉路口仿真中，本文方法在CFCR（$68.4\%$ vs.\ MAPPO $57.8\%$，$p<0.001$，$d=5.27$）和MSVR（$2.4\%$ vs.\ $3.0\%$，$p=0.029$，$d=1.69$）上取得统计显著的改善；TC的差异（$18.2$ vs.\ $20.3$，降幅10.3\%）有大效应量（$d=0.93$）但受$n=5$限制未达显著，有待$n\geq20$的重复实验确认。消融实验表明DyGAT模块是贡献最大的组件。

从实践角度看，本文结果对低空交通管理系统（UTM）的规划设计具有两点启示。其一，安全评估指标的选择至关重要：以CFCR为代表的“零风险接触”类指标，应作为低空交通安全性能评价的核心口径，单纯依赖通过率会系统性高估系统的安全性。其二，物理先验与学习机制的融合是一条可行路径：本文以约0.2 ms的额外推理开销（向量化实现）换取可度量的安全提升，即可以在算力与安全裕度之间以领域知识换取安全收益，这一特性对资源受限的机载计算环境具有实际价值。

后续工作分为两组。**其一，支撑本文结论的核心实验（优先级最高）**：（1）DyGAT三路组件级消融（仅$w_{ij}$/仅$e_{ij}$/$e_{ij}\odot w_{ij}$），验证“学习与门控缺一不可”的核心假设；（2）按$r$分级统计各严重等级分布，并将严重冲突率对标民航安全目标；（3）$n\geq20$重复实验，以配对/置换检验与bootstrap复核本文的统计结论。**其二，模型的扩展方向**：（4）三维多路口网络与高度层动态分配；（5）集成Dryden/Kaimal标准风谱模型与通信延迟/丢包约束，检验非理想环境下的性能退化；（6）GAT-Mixer与均值池化/QMIX混合网络的直接对比，以及与MPC/CBF的混合架构（“学习决策+安全兜底”）；（7）在条件成熟时将模型迁移至硬件在环仿真平台验证。

\textbf{代码与可复现性声明：}实验代码（含AirTrafficEnv仿真环境、DyGAT+GAT-Mixer模型实现、ORCA基线适配、训练与评估脚本）将在论文接收后于GitHub开源（匿名链接在camera-ready版本提供）。训练计算量方面，本文报告以实测吞吐为参照（RTX 4060 GPU）：训练单步平均耗时约0.1--0.3 s/步（含LSTM/DyGAT前向、策略与价值网络的前向与反向传播及4轮PPO更新摊薄；$N=5$时约9 s/episode，随$N$与episode长度增大而上升）。由于正文未公开确切的总训练步数口径，此处不给出固定的“单seed小时数”，作者将在代码开源时一并公布实际训练预算与实测吞吐，以保证复现成本可核验。第5.9节的推理单步耗时（$N=21$时约213 ms）与训练单步成本是不同口径，不应混用。**可复现性说明：** 正文所报告的实验结果（第5章各表）对应完整训练配置（含课程学习全程、$\lambda=2.0$冲突加权及完整训练预算），该配置与训练脚本将在代码发布时一并提供；本文同时公开各补充实验的可运行脚本（\textit{experiment\_scripts/} 目录）及其实测吞吐/推理数据，供第三方按相同口径核验。

%==========================================================================
% 参考文献
%==========================================================================
\begin{thebibliography}{99}
\bibitem{ref45} COHEN A P, SHAHEEN S A, FARRAR E M. Urban air mobility: history, ecosystem, market potential, and challenges[J]. IEEE Transactions on Intelligent Transportation Systems, 2021, 22(9): 6074--6087.

\bibitem{ref1} Federal Aviation Administration. Urban air mobility (UAM): concept of operations v2.0[R]. Washington, DC: FAA, 2023.

\bibitem{ref46} 国务院中央军委空中交通管制委员会, 中国民用航空局. 国家空域基础分类方法[S]. 北京, 2023.

\bibitem{ref48} International Civil Aviation Organization. Global air navigation plan (GANP) 2016--2030: Doc 9750-AN/963[R]. 5th ed. Montreal: ICAO, 2016.

\bibitem{ref47} 中国民用航空局. 民用无人驾驶航空器运行安全管理规则（CCAR-92）[S]. 北京: 中国民用航空局, 2023.

\bibitem{ref49} International Civil Aviation Organization. Air traffic management: Doc 4444[R]. 16th ed. Montreal: ICAO, 2016.

\bibitem{ref2} MENON P K, SWERIDUK G D, BILIMORIA K D. New approach to modeling, analysis, and control of air traffic flow[J]. Journal of Guidance, Control, and Dynamics, 2004, 27(5): 737--744.

\bibitem{ref12} ASLANI M, MESGARI M S, WIERING M. Adaptive traffic signal control with actor-critic methods in a real-world traffic network with different traffic disruption events[J]. Transportation Research Part C: Emerging Technologies, 2017, 85: 732--752.

\bibitem{ref16} 赵晓华, 石建军, 李振龙, 赵国勇. 基于Q-learning和BP神经元网络的交叉口信号灯控制[J]. 公路交通科技, 2007, 24(7): 99--102.

\bibitem{ref19} 承向军, 常歆识, 杨肇夏. 基于Q-学习的交通信号控制方法[J]. 系统工程理论与实践, 2006(8): 136--140.

\bibitem{ref20} 卢守峰, 张术, 刘喜敏. 平均排队长度差最小的单交叉口在线Q学习模型[J]. 公路交通科技, 2014, 31(11): 116--122.

\bibitem{ref15} 王迎. 城市交通信号控制深度强化学习算法研究[D]. 济南: 山东交通学院, 2024.

\bibitem{ref17} 赖建辉. 基于深度强化学习的交通控制优化方法研究及实现[D]. 杭州: 浙江工业大学, 2020.

\bibitem{ref18} 张新武. 基于深度强化学习算法的交通信号控制优化研究[D]. 太原: 太原科技大学, 2024.

\bibitem{ref22} 张春阳, 席雅琪. 人工智能在城市交通控制中的应用综述[J]. 信息系统工程, 2024(07): 33--36.

\bibitem{ref3} RICHARDS A, HOW J P. Aircraft trajectory planning with collision avoidance using mixed integer linear programming[C]//Proceedings of the 2002 American Control Conference. Piscataway: IEEE, 2002: 1936--1941.

\bibitem{ref4} VAN DEN BERG J, GUY S J, LIN M, et al. Reciprocal n-body collision avoidance[C]//Robotics Research: The 14th International Symposium ISRR (Springer Tracts in Advanced Robotics, Vol.\ 70). Berlin: Springer, 2011: 3--19.

\bibitem{ref5} AMES A D, COOGAN S, EGERSTEDT M, et al. Control barrier functions: theory and applications[C]//2019 18th European Control Conference (ECC). Piscataway: IEEE, 2019: 3420--3431.

\bibitem{ref40} MNIH V, KAVUKCUOGLU K, SILVER D, et al. Human-level control through deep reinforcement learning[J]. Nature, 2015, 518(7540): 529--533.

\bibitem{ref21} 江波, 李彦冬, 刘威陇, 等. 应用深度Q学习的机场道口协同智能调速方法[J]. 航空计算技术, 2022, 52(1): 26--30.

\bibitem{ref6} AGOGINO A, TUMER K. Regulating air traffic flow with coupled agents[C]//Proceedings of the 7th International Conference on Autonomous Agents and Multiagent Systems (AAMAS 2008). Richland: IFAAMAS, 2008: 535--542.

\bibitem{ref7} AGOGINO A K, TUMER K. A multiagent approach to managing air traffic flow[J]. Autonomous Agents and Multi-Agent Systems, 2012, 24(1): 1--25.

\bibitem{ref14} TUMER K, AGOGINO A. Distributed agent-based air traffic flow management[C]//Proceedings of the 6th International Joint Conference on Autonomous Agents and Multiagent Systems (AAMAS 2007). Honolulu: ACM, 2007: 330--337.

\bibitem{ref24} TUMER K, AGOGINO A. Improving air traffic management with a learning multiagent system[J]. IEEE Intelligent Systems, 2009, 24(1): 18--21.

\bibitem{ref11} SPATHARIS C, BASTAS A, KRAVARIS T, et al. Hierarchical multiagent reinforcement learning schemes for air traffic management[J]. Neural Computing and Applications, 2021, 33(14): 8649--8671.

\bibitem{ref23} GHAVAMZADEH M, MAHADEVAN S, MAKAR R. Hierarchical multi-agent reinforcement learning[J]. Autonomous Agents and Multi-Agent Systems, 2006, 13(2): 197--229.

\bibitem{ref8} BRITTAIN M, WEI P. Autonomous separation assurance in a high-density en route sector: a deep multi-agent reinforcement learning approach[C]//2019 IEEE Intelligent Transportation Systems Conference (ITSC). Piscataway: IEEE, 2019: 3256--3262.

\bibitem{ref9} GHOSH S, LAGUNA S, LIM S H, et al. A deep ensemble method for multi-agent reinforcement learning: a case study on air traffic control[C]//Proceedings of the 31st International Conference on Automated Planning and Scheduling (ICAPS 2021). Palo Alto: AAAI Press, 2021: 468--476.

\bibitem{ref29} XIE Y B, GARDI A, SABATINI R. Reinforcement learning-based flow management techniques for urban air mobility and dense low-altitude air traffic operations[C]//2021 IEEE/AIAA 40th Digital Avionics Systems Conference (DASC). Piscataway: IEEE, 2021: 1--10.

\bibitem{ref34} LOWE R, WU Y, TAMAR A, et al. Multi-agent actor-critic for mixed cooperative-competitive environments[C]//Advances in Neural Information Processing Systems 30 (NeurIPS 2017). Long Beach: Curran Associates, 2017: 6379--6390.

\bibitem{ref35} FOERSTER J, FARQUHAR G, AFOURAS T, et al. Counterfactual multi-agent policy gradients[C]//Proceedings of the 32nd AAAI Conference on Artificial Intelligence (AAAI 2018). New Orleans: AAAI Press, 2018: 2974--2982.

\bibitem{ref33} RASHID T, SAMVELYAN M, DE WITT C S, et al. QMIX: monotonic value function factorisation for deep multi-agent reinforcement learning[C]//Proceedings of the 35th International Conference on Machine Learning (ICML 2018). Stockholmsm\“{a}ssan: PMLR, 2018: 4295--4304.

\bibitem{ref36} YU C, VELU A, VINITSKY E, et al. The surprising effectiveness of PPO in cooperative multi-agent games[C]//Advances in Neural Information Processing Systems 35 (NeurIPS 2022). New Orleans: Curran Associates, 2022: 24611--24624.

\bibitem{ref13} 翟子洋, 郝茹茹, 董世浩. 大规模智慧交通信号控制中的强化学习和深度强化学习方法综述[J]. 计算机应用研究, 2024, 41(6): 1618--1627.

\bibitem{ref32} KIPF T N, WELLING M. Semi-supervised classification with graph convolutional networks[C]//Proceedings of the 5th International Conference on Learning Representations (ICLR 2017). Toulon: ICLR, 2017.

\bibitem{ref31} VELI\v{C}KOVI\'C P, CUCURULL G, CASANOVA A, et al. Graph attention networks[C]//Proceedings of the 6th International Conference on Learning Representations (ICLR 2018). Vancouver: ICLR, 2018.

\bibitem{ref37} HOCHREITER S, SCHMIDHUBER J. Long short-term memory[J]. Neural Computation, 1997, 9(8): 1735--1780.

\bibitem{ref55} KARNIADAKIS G E, KEVREKIDIS I G, LU L, et al. Physics-informed machine learning[J]. Nature Reviews Physics, 2021, 3(6): 422--440.

\bibitem{ref10} LI Y, LI J, WANG J, ZHANG X, DING H, DU W. Multiscale-graph-enhanced reinforcement learning for conflict resolution in dense UAV networks[J]. IEEE Internet of Things Journal, 2025, 12(21): 44290--44303.

\bibitem{ref41} BERNSTEIN D S, GIVAN R, IMMERMAN N, et al. The complexity of decentralized control of Markov decision processes[J]. Mathematics of Operations Research, 2002, 27(4): 819--840.

\bibitem{ref25} BACCHINI A, CESTINO E. Electric VTOL configurations comparison[J]. Aerospace, 2019, 6(3): 26.

\bibitem{ref26} U.S. Department of Defense. Flying qualities of piloted airplanes: MIL-F-8785C[S]. Washington, DC: U.S. Department of Defense, 1980.

\bibitem{ref44} International Electrotechnical Commission. Wind energy generation systems -- Part 1: design requirements: IEC 61400-1[S]. 4th ed. Geneva: IEC, 2019.

\bibitem{ref51} RTCA. DO-365: Minimum operational performance standards (MOPS) for detect and avoid systems for unmanned aircraft systems[S]. Washington, DC: RTCA, 2017.

\bibitem{ref50} Joint Authorities for Rulemaking of Unmanned Systems. JARUS guidelines on specific operations risk assessment (SORA) v2.0[S]. JARUS, 2019.

\bibitem{ref28} CONG W, ZHANG S, CHEN J, et al. Do we really need complicated model architectures for temporal networks?[C]//Proceedings of the 11th International Conference on Learning Representations (ICLR 2023). Kigali: ICLR, 2023.

\bibitem{ref39} BENGIO Y, LOUGARD J, COLLOBERT R, et al. Curriculum learning[C]//Proceedings of the 26th International Conference on Machine Learning (ICML 2009). Montreal: ACM, 2009: 41--48.

\bibitem{ref30} SCHULMAN J, WOLSKI F, DHARIWAL P, et al. Proximal policy optimization algorithms[EB/OL]. arXiv:1707.06347, 2017.

\bibitem{ref27} LIANG X, MA Y, FENG Y, et al. PTR-PPO: proximal policy optimization with prioritized trajectory replay[EB/OL]. arXiv:2112.03798, 2021.

\bibitem{ref38} SCHULMAN J, MORITZ P, LEVINE S, et al. High-dimensional continuous control using generalized advantage estimation[C]//Proceedings of the 4th International Conference on Learning Representations (ICLR 2016). San Juan: ICLR, 2016.

\bibitem{ref52} HOEKSTRA J M, ELLEBROEK J. Bluesky ATC simulator project: an open data and open source approach[C]//Proceedings of the 7th International Conference on Research in Air Transportation (ICRAT 2016). 2016.

\bibitem{ref53} SCH\”{A}FER M, STROHMEIER M, LENDERS V, et al. Bringing up OpenSky: a large-scale ADS-B sensor network for research[C]//Proceedings of the 13th International Symposium on Information Processing in Sensor Networks (IPSN). Piscataway: IEEE, 2014: 83--94.

\bibitem{ref54} LOPEZ P A, BEHRISCH M, BIEKER-WALZ L, et al. Microscopic traffic simulation using SUMO[C]//2018 21st International Conference on Intelligent Transportation Systems (ITSC). Piscataway: IEEE, 2018: 2575--2582.

\bibitem{ref42} WELCH B L. The generalization of Student's problem when several different population variances are involved[J]. Biometrika, 1947, 34(1--2): 28--35.

\bibitem{ref43} COHEN J. Statistical power analysis for the behavioral sciences[M]. 2nd ed. Hillsdale: Lawrence Erlbaum Associates, 1988.
\end{thebibliography}

\end{document}
